% >>> GENERALIZATION (Fable 5.1, Sep 1 2026): the audit gains a paragraph showing the sorting is the coupling lemma +
% Uhlmann + Fuchs-van de Graaf; five bib keys added VERIFY-flagged.
% >>> PEDAGOGY PASS 2 (Fable 5.1, Sep 1 2026): the 'exclusive regime enters' transition moved to follow the
% record sentence it announces; the 887-word walk split into three paragraphs; the two-relations
% clarification (TV split vs square map) and the 'unconditional' label carried here from the Brief.
% >>> PROPAGATION (Fable 5, Sep 1 2026): the stake sentence reframed off 'record-complexity ceiling' / 'names a rung' onto
% the physics of records (the no-cut position). See the ledger, R123.
% >>> PEDAGOGY (Fable 5, Aug 29 2026): the V/D roadmap at \S twotag (medium treatment, framework-neutral), the
% displayed port identity (eq:portidentity), and the per-result audit at the section's end. See the ledger, R113.
% >>> PROPAGATION (Fable 5, Aug 29 2026): the negativity-compliance sentence aligned with the signed-weight doctrine
% (compliance rides on the single-sample-space refusal); the extraction-axis phrase disambiguated (orthogonal to the
% markers' bisector on the sphere, the 45-degree straddle, not the erasing bisector); the continuum extension of the
% classification added to the frontier list. See the ledger, R109.
% ============================================================================
% EXTRACTION PAPER 1 (physics-style): the reconstruction-genre presentation.
% Drafted by Claude Fable 5, July 2026, per the author's commission, as a
% compression/re-registration of the author's short paper (itself the
% compression of the long paper). All technical content is the author's;
% the genre framing, placement prose, and compression choices are the
% GRADING PASS (Fable 5, August 29, 2026), propagated from the Brief's referee
% rounds. Fenced "% >>> EDIT Qn". (Q1) RQM comparison regraded in both places:
% the relativization is postulated as RQM postulates it, and what is derived is
% the armor. (Q2) "prediction in kind" --> "commitment to a prediction", with
% the echo-regime and ladder-from-below disciplines added at the costs section.
% (Q3) Theorem setup indexed per Postulate 1, with the two delayed-choice
% families sorted by the regime criterion (Wheeler vs quantum-control).
% (Q4) "Collapse is derived" graded: form from the elimination, trigger from
% the regime identification. (Q5) No-signaling marked a placement result.
% New bib entry ionicioiu2011proposal --- VERIFY at submission.
% drafting AI's, for the author's revision. VERIFY-flagged bib entries
% DELTA AUDIT (July 2026, added at the verification round's direction): the two
% sentences above described the original compression and are INACCURATE for this
% draft's delta. Content beyond the deposited version --- the triangle
% characterization, the monotone re-classification theorem, the leakage-geometry
% analysis of the Beyond-the-Circle section, and the compressed conditional route
% (a relocated standard discrimination argument) --- is of mixed provenance: built
% from the author's postulates, framework, and conversational contributions (the
% region-node dial, the sequencing constraint, the two-tag conception), with the
% derivation drafting, standard-formalism translations, and same-family referee
% corrections supplied by the drafting model. A dated section-level audit travels
% in the companion session record. Wording subject to the author's revision;
% independent human review is the gate on this item specifically.
% should be checked against the sources before submission.
% ============================================================================
\begin{filecontents}[overwrite]{recon_refs.bib}
% referee's read --- add it there, whereupon (refs.bib loading first) this local becomes
% a shadowed duplicate and may be removed.
@comment{EDIT Q3 (Fable 5, Aug 29 2026): added for the delayed-choice sorting; fields from working knowledge --- VERIFY at submission.}
@article{ionicioiu2011proposal, author={Ionicioiu, Radu and Terno, Daniel R.}, title={Proposal for a quantum delayed-choice experiment}, journal={Physical Review Letters}, volume={107}, number={23}, pages={230406}, year={2011}}
@misc{schmid2020unscrambling, author={Schmid, David and Selby, John H. and Spekkens, Robert W.}, title={Unscrambling the omelette of causation and inference: The framework of causal-inferential theories}, year={2020}, note={arXiv:2009.03297}}
@comment{EDIT B1 (Fable 5, Aug 29 2026): schmid2020unscrambling --- subtitle per referee recollection, VERIFY before submission. Formerly part of the visible note field.}
% referee's read and should be added there at merge.
@article{hardy2001,
  author = {Hardy, Lucien},
  title = {Quantum Theory From Five Reasonable Axioms},
  journal = {arXiv preprint},
  eprint = {quant-ph/0101012},
  year = {2001}
}
@article{cbh2003,
  author = {Clifton, Rob and Bub, Jeffrey and Halvorson, Hans},
  title = {Characterizing Quantum Theory in Terms of Information-Theoretic Constraints},
  journal = {Foundations of Physics},
  volume = {33}, pages = {1561--1591}, year = {2003}
}
@article{cdp2011,
  author = {Chiribella, Giulio and D'Ariano, Giacomo Mauro and Perinotti, Paolo},
  title = {Informational derivation of quantum theory},
  journal = {Physical Review A},
  volume = {84}, pages = {012311}, year = {2011}
}
@article{masanesmuller2011,
  author = {Masanes, Llu{\'i}s and M{\"u}ller, Markus P.},
  title = {A derivation of quantum theory from physical requirements},
  journal = {New Journal of Physics},
  volume = {13}, pages = {063001}, year = {2011}
}
@article{spekkens2007,
  author = {Spekkens, Robert W.},
  title = {Evidence for the epistemic view of quantum states: A toy theory},
  journal = {Physical Review A},
  volume = {75}, pages = {032110}, year = {2007}
}
@article{rovelli1996,
  author = {Rovelli, Carlo},
  title = {Relational quantum mechanics},
  journal = {International Journal of Theoretical Physics},
  volume = {35}, pages = {1637--1678}, year = {1996}
}
@article{zurek2003,
  author = {Zurek, Wojciech H.},
  title = {Environment-Assisted Invariance, Entanglement, and Probabilities in Quantum Physics},
  journal = {Physical Review Letters},
  volume = {90}, pages = {120404}, year = {2003}
}
% VERIFY: Dakic-Brukner venue/pages before submission.
@incollection{dakicbrukner2011,
  author = {Daki{\'c}, Borivoje and Brukner, {\v C}aslav},
  title = {Quantum Theory and Beyond: Is Entanglement Special?},
  booktitle = {Deep Beauty: Understanding the Quantum World through Mathematical Innovation},
  editor = {Halvorson, Hans},
  publisher = {Cambridge University Press},
  pages = {365--392},
  year = {2011}
}
@article{griffiths1984,
  author = {Griffiths, Robert B.},
  title = {Consistent histories and the interpretation of quantum mechanics},
  journal = {Journal of Statistical Physics},
  volume = {36}, pages = {219--272}, year = {1984}
}
@incollection{gellmannhartle1990,
  author = {Gell-Mann, Murray and Hartle, James B.},
  title = {Quantum Mechanics in the Light of Quantum Cosmology},
  booktitle = {Complexity, Entropy and the Physics of Information},
  publisher = {Addison-Wesley},
  year = {1990}
}
@article{bell1964,
  author = {Bell, John S.},
  title = {On the Einstein Podolsky Rosen paradox},
  journal = {Physics Physique Fizika},
  volume = {1}, pages = {195--200}, year = {1964}
}
@article{epr1935,
  author = {Einstein, Albert and Podolsky, Boris and Rosen, Nathan},
  title = {Can Quantum-Mechanical Description of Physical Reality Be Considered Complete?},
  journal = {Physical Review},
  volume = {47}, pages = {777--780}, year = {1935}
}
@article{fine1982,
  author = {Fine, Arthur},
  title = {Hidden Variables, Joint Probability, and the Bell Inequalities},
  journal = {Physical Review Letters},
  volume = {48}, pages = {291--295}, year = {1982}
}
@article{pitowsky1994,
  author = {Pitowsky, Itamar},
  title = {George Boole's `Conditions of Possible Experience' and the Quantum Puzzle},
  journal = {The British Journal for the Philosophy of Science},
  volume = {45}, pages = {95--125}, year = {1994}
}
@book{feynman1965lectures,
  author = {Feynman, Richard P. and Leighton, Robert B. and Sands, Matthew},
  title = {The Feynman Lectures on Physics, Vol.~III: Quantum Mechanics},
  publisher = {Addison-Wesley},
  address = {Reading, MA},
  year = {1965}
}
@article{feynman1948,
  author = {Feynman, Richard P.},
  title = {Space-Time Approach to Non-Relativistic Quantum Mechanics},
  journal = {Reviews of Modern Physics},
  volume = {20}, pages = {367--387}, year = {1948}
}
@article{dirac1933,
  author = {Dirac, Paul A. M.},
  title = {The Lagrangian in Quantum Mechanics},
  journal = {Physikalische Zeitschrift der Sowjetunion},
  volume = {3}, pages = {64--72}, year = {1933}
}
@article{greenberger1988,
  author = {Greenberger, Daniel M. and Yasin, Allaine},
  title = {Simultaneous wave and particle knowledge in a neutron interferometer},
  journal = {Physics Letters A},
  volume = {128}, pages = {391--394}, year = {1988}
}
@article{englert1996,
  author = {Englert, Berthold-Georg},
  title = {Fringe Visibility and Which-Way Information: An Inequality},
  journal = {Physical Review Letters},
  volume = {77}, pages = {2154--2157}, year = {1996}
}
@article{scully1991,
  author = {Scully, Marlan O. and Englert, Berthold-Georg and Walther, Herbert},
  title = {Quantum optical tests of complementarity},
  journal = {Nature},
  volume = {351}, pages = {111--116}, year = {1991}
}
@comment{VERIFY at submission (entered from working knowledge): lindvall1992lectures, uhlmann1976transition, jozsa1994fidelity, fuchs1999cryptographic, durr2001quantitative --- added Sep 1 2026 for the generalization paragraph.}
@book{lindvall1992lectures, author={Lindvall, Torgny}, title={Lectures on the Coupling Method}, publisher={Wiley}, address={New York}, year={1992}}
@article{uhlmann1976transition, author={Uhlmann, Armin}, title={The ``transition probability'' in the state space of a $*$-algebra}, journal={Reports on Mathematical Physics}, volume={9}, number={2}, pages={273--279}, year={1976}}
@article{jozsa1994fidelity, author={Jozsa, Richard}, title={Fidelity for mixed quantum states}, journal={Journal of Modern Optics}, volume={41}, number={12}, pages={2315--2323}, year={1994}}
@article{fuchs1999cryptographic, author={Fuchs, Christopher A. and van de Graaf, Jeroen}, title={Cryptographic distinguishability measures for quantum-mechanical states}, journal={IEEE Transactions on Information Theory}, volume={45}, number={4}, pages={1216--1227}, year={1999}}
@article{durr2001quantitative, author={D{\"u}rr, Stephan}, title={Quantitative wave-particle duality in multibeam interferometers}, journal={Physical Review A}, volume={64}, pages={042113}, year={2001}}
@article{durr1998,
  author = {D{\"u}rr, S. and Nonn, T. and Rempe, G.},
  title = {Origin of quantum-mechanical complementarity probed by a `which-way' experiment in an atom interferometer},
  journal = {Nature},
  volume = {395}, pages = {33--37}, year = {1998}
}
@article{arndt1999,
  author = {Arndt, Markus and others},
  title = {Wave--particle duality of C$_{60}$ molecules},
  journal = {Nature},
  volume = {401}, pages = {680--682}, year = {1999}
}
@article{fein2019,
  author = {Fein, Yaakov Y. and others},
  title = {Quantum superposition of molecules beyond 25 kDa},
  journal = {Nature Physics},
  volume = {15}, pages = {1242--1245}, year = {2019}
}
@incollection{wheeler1983,
  author = {Wheeler, John A.},
  title = {Law without law},
  booktitle = {Quantum Theory and Measurement},
  publisher = {Princeton University Press},
  year = {1983}
}
@article{peres1978,
  author = {Peres, Asher},
  title = {Unperformed experiments have no results},
  journal = {American Journal of Physics},
  volume = {46}, pages = {745--747}, year = {1978}
}
@article{frauchiger2018,
  author = {Frauchiger, Daniela and Renner, Renato},
  title = {Quantum theory cannot consistently describe the use of itself},
  journal = {Nature Communications},
  volume = {9}, pages = {3711}, year = {2018}
}
@article{pusey2012,
  author = {Pusey, Matthew F. and Barrett, Jonathan and Rudolph, Terry},
  title = {On the reality of the quantum state},
  journal = {Nature Physics},
  volume = {8}, pages = {475--478}, year = {2012}
}
@misc{AethusTree, author={Benander, Ajax}, title={The Weight Algebra of Aethae: A Two-Layer Commutative Semiring for Superposition and Pruning}, year={2026}, howpublished={Preprint}, url={https://philpapers.org/rec/BENTWA-5}}
@misc{AethusUnion, author={Benander, Ajax}, title={Exclusivity as Achievement: Relocating Mutual Exclusion to the Interior, and the Composition Regime That Falls Out}, year={2026}, howpublished={Preprint}, url={https://philpapers.org/rec/BENEAA-9}}
@misc{AethusBrief, author={Benander, Ajax}, title={Aethic Reasoning in Brief: From a Counterfactual Semantics to the Quantum Measurement Problem}, year={2026}, howpublished={Preprint}, url={https://philpapers.org/rec/BENARI-3}, eprint={https://philsci-archive.pitt.edu/id/eprint/30439}, eprinttype={Alternative URL}}
@article{jakob2010quantitative,
  author = {Jakob, Matthias and Bergou, J{\'a}nos A.},
  title = {Quantitative complementarity relations in bipartite systems: Entanglement as a physical reality},
  journal = {Optics Communications},
  volume = {283}, number = {5}, pages = {827--830}, year = {2010},
  doi = {10.1016/j.optcom.2009.10.044}
}
@article{wood_spekkens_2015,
  author = {Wood, Christopher J. and Spekkens, Robert W.},
  title = {The lesson of causal discovery algorithms for quantum correlations: 
           Causal explanations of Bell-inequality violations require fine-tuning},
  journal = {New Journal of Physics},
  volume = {17},
  number = {3},
  pages = {033002},
  year = {2015}
}
@article{berry1984quantal,
  author  = {Berry, Michael V.},
  title   = {Quantal phase factors accompanying adiabatic changes},
  journal = {Proceedings of the Royal Society of London A},
  volume  = {392},
  number  = {1802},
  pages   = {45--57},
  year    = {1984}
}
 
@article{Kochetov1995,
  author  = {Kochetov, E. A.},
  title   = {{SU(2)} coherent-state path integral},
  journal = {Journal of Mathematical Physics},
  volume  = {36},
  number  = {9},
  pages   = {4667--4679},
  year    = {1995},
  doi     = {10.1063/1.530913}
}
 
@article{pancharatnam1956generalized,
  author  = {Pancharatnam, S.},
  title   = {Generalized theory of interference, and its applications. {Part I}: {Coherent} pencils},
  journal = {Proceedings of the Indian Academy of Sciences, Section A},
  volume  = {44},
  number  = {5},
  pages   = {247--262},
  year    = {1956},
  doi     = {10.1007/BF03046050}
}

@misc{Aethus, author = {Benander, Ajax}, title = {Aethic Reasoning: A Comprehensive Solution to the Quantum Measurement Problem}, howpublished = {Preprint}, year = {2026}, url = {https://philpapers.org/rec/BENARA-5}, note = {The full development; PhilPapers record BENARA-5}}

@article{Joos1985, author={Joos, E. and Zeh, H. D.}, title={The emergence of classical properties through interaction with the environment}, journal={Zeitschrift f\"ur Physik B}, volume={59}, pages={223--243}, year={1985}}

@incollection{jaynes1990probability, author={Jaynes, E. T.}, title={Probability in Quantum Theory}, booktitle={Complexity, Entropy, and the Physics of Information}, editor={Zurek, W. H.}, publisher={Addison-Wesley}, year={1990}}

@misc{spekkens2018quantum,
    author = "Spekkens, Robert",
    title = "{Robert Spekkens Public Lecture: The Riddle of the Quantum Sphinx}",
    howpublished = "YouTube",
    year = "2018",
    month = "February",
    day = "8",
    note = "Perimeter Institute Public Lecture Series",
    url = "https://www.youtube.com/watch?v=jvZ5m3UMer4"
}

@misc{spekkens2024escapedsapiens,
  author       = {Spekkens, Robert W.},
  title        = {What is Reality Made Of? {Unscrambling} the Quantum Omelet},
  howpublished = {Interview by Shane Farnsworth, The Escaped Sapiens Podcast, Episode 74},
  year         = {2024},
  month        = nov,
  note         = {Video available at \url{https://youtu.be/J2ZIRkfrFlI}},
  url          = {https://youtu.be/J2ZIRkfrFlI}
}

@misc{GoldenScience2022,
  author       = {Benander, Ajax},
  title        = {Golden Science},
  year         = {2022},
  month        = {November},
  howpublished = {Self-hosted manuscript},
  url          = {https://aethus.org/aethic-code/files/GoldenScience_Nov2022.pdf},
  note         = {Title page August 2022, the date of the first version of the source; this version created November 2022. Not deposited to a
                  preprint server; cited as a dated developmental record rather than
                  as a source of results.}
}

@article{elitzur1993quantum,
  author  = {Elitzur, Avshalom C. and Vaidman, Lev},
  title   = {Quantum Mechanical Interaction-Free Measurements},
  journal = {Foundations of Physics},
  volume  = {23},
  number  = {7},
  pages   = {987--997},
  year    = {1993},
  doi     = {10.1007/BF00736012}
}
\end{filecontents}

\documentclass[11pt]{article}
\usepackage[margin=0.95in]{geometry}
\usepackage{amsmath, amssymb, amsthm}
\usepackage[hidelinks]{hyperref}
\usepackage{microtype}
\usepackage{booktabs}
\usepackage[style=numeric, sorting=none]{biblatex}
\addbibresource{refs.bib}
\addbibresource{recon_refs.bib}
\newtheorem{theorem}{Theorem}
\newtheorem{proposition}{Proposition}
\newtheorem{corollary}{Corollary}
\newtheorem{lemma}{Lemma}
\newtheorem{claim}[lemma]{Claim}
\theoremstyle{definition}
\newtheorem{postulate}{Postulate}

\title{Quantum Mechanics from Records:\\ \large Measurement Structure on a Finite Path-Kernel Domain}
\author{Ajax Benander\thanks{Email: \href{mailto:akb9408@rit.edu}{akb9408@rit.edu}, Website: \url{https://aethus.org}. This paper is a self-contained extraction from a larger framework; the companion papers \cite{AethusBrief, Aethus} develop the full structure, including the modal semantics from which the present postulates were originally abstracted.}}
\date{First posted July 17, 2026\\ {\small This version: September 2026 --- draft}}

\begin{document}
\maketitle

\begin{abstract}
\noindent We give a finite path-kernel representation of record-conditioned quantum statistics, with a qualitative admissibility classification and an explicit restriction operation. The primitives are consistency constraints on observer-indexed partial information states: determinacy relative to a state, unanswered attributes read as superposition, and a pruning rule under which a state is valid when some declared refinement of it has no contradiction-reaching descendant. Pruning ranges over admissible, non-negligible components at every refinement depth, record cells among them. Under a complete and persistent record structure, where the record signatures include every relevant accessible record and every declared refinement retains a separating record, we prove that this rule classifies every subset of a space of paths as valid, over-specified, or consistency-pruned, so that, under the stated visibility assumptions, the quantum--classical boundary is the record partition, determined before the quantitative path-kernel weights are assigned. One quantitative axiom, the path kernel, then yields on a fixed finite domain a Hilbert-space representation of the admissible clusters with the Born statistics, record observables as self-adjoint operators, and the projection postulate as the representation of a stipulated record-compatible restriction, realized as L\"{u}ders conditioning and not selected by the validity rule alone; the representation does not reach every Hilbert state on any one finite domain. The two-tag interferometer configuration exhibits both composition regimes in one apparatus and meets the quantitative predictions, and the same admissibility criterion applies to an unread which-path record and to an unread friend's laboratory, where the record premise for a whole laboratory remains a physical hypothesis.
\end{abstract}

\section{Introduction}\label{sec:intro}

The Dirac--von Neumann axioms present quantum mechanics with the state as primitive and measurement as a separate rule bolted onto the dynamics. This paper reverses the order of presentation on one finite domain and reports exactly what the reversal buys. The structure reconstructed is the measurement architecture of the standard formalism: Hilbert representation of cluster states, Born statistics, record observables, composite systems, unitary evolution, and the projection postulate. The inputs are stated in one package. Three postulates on observer-indexed information states fix determinacy as a two-place relation, identify informational blanks with superpositions, and prune candidate states by a two-quantifier safety rule, the third postulate; a record axiom says what a separating record is; a composition clause distinguishes alternatives that exclude one another from alternatives that do not; one quantitative axiom, the path kernel, supplies weights; a bridge principle carries formal invalidity to physical unrealizability; and one stipulated record-compatible restriction supplies the update. What is proved: a qualitative classification theorem, that under a complete and persistent record structure every subset of a space of paths is valid, over-specified or consistency-pruned, with the quantum--classical boundary the record partition; a representation theorem, that given the kernel the classification's clusters are represented in a Hilbert space with the Born statistics, record observables as self-adjoint operators by a soundness result, and the stipulated restriction as L\"{u}ders conditioning, $\Phi(T \cap E_j) = P_j \Phi(T)$; a local polarization lemma, recovering amplitudes from weights at each configuration; and a no-signaling identity for a fixed record-signature partition. Completeness of the state and observable representations is excluded on any one fixed finite domain, and the physical realization of the bridge is an open construction, named as such.

What makes the reformulation illuminating is where it places the measurement problem. The projection postulate is not eliminated but relocated: the validity postulate alone does not select the update, and what the representation theorem shows is that a stipulated record-compatible restriction, once supplied, is realized by the cluster representation as projection. Under the stated structural-visibility assumptions, the qualitative classification precedes the assignment of quantitative path-kernel weights: the boundary is drawn by which records the arrangement makes available and which of them persist through its refinements, and the quantitative half is the kernel's. Two consequences are worked in full: the two-tag configuration, in which both composition regimes stand in one apparatus and the classification meets the quantitative predictions; and the detector--friend thesis, that an unread which-path record and an unread friend's outcome record fall under one admissibility criterion, with the condition that must still be established for a whole laboratory stated as a physical premise.

\begin{table}[htbp]
\centering
\small
\caption{Dependency map: each result with the inputs it consumes. The first two rows use no quantitative axiom; the next eight are conditional on the kernel $K$; the last needs a further physical premise.}
\label{tab:kernel-boundary}
\begin{tabular}{p{0.27\textwidth} p{0.45\textwidth} p{0.2\textwidth}}
\toprule
\textbf{Result} & \textbf{Inputs that must be explicit} & \textbf{Where} \\
\midrule
Record-partition classification & Declared descent; a complete and persistent record structure; every record cell structurally visible; record-individuation for the over-specification case & Theorem~\ref{thm:powerset} \\
The two composition regimes; the chord $V + D = 1$ & The regimes of the second postulate; convex weights, with ideal endpoint patterns, compatible normalization and an accessible tag for the chord & Section~\ref{sec:regimes} \\
State representation & The finite path domain and the kernel (K1); $W(T) > 0$; one ray per such cluster, the representation not injective & Theorem~\ref{thm:recovery}~(i) \\
Superposition as vector addition & (K1); disjoint clusters, coefficients $(1,1)$ & Theorem~\ref{thm:recovery}~(ii) \\
Probability assignment & The weight law (K1), the calibration (K3), normalization & Theorem~\ref{thm:recovery}~(iv) \\
Record observables & The record partition, a record map factoring through the endpoints, the representation; soundness only & Theorem~\ref{thm:recovery}~(iii) \\
Projective conditioning & The stipulated record-compatible restriction and its representation; an outcome of positive weight & Theorem~\ref{thm:recovery}~(v) \\
Composites and evolution & The product slice; unitarity (K2) & Theorem~\ref{thm:recovery}~(vi)--(vii) \\
Amplitudes from weights & (K1); at each configuration, up to a common phase and conjugation; the global reconstruction open & Lemma~\ref{lem:polarization} \\
No-signaling & The two composition laws on a fixed signature partition & Theorem~\ref{thm:recovery}~(viii) \\
Laboratory application & An additional physical record premise & Lemma~\ref{lem:record-exclusion-q} \\
\bottomrule
\end{tabular}
\end{table}

\noindent The dependency structure, stated once so that it need not be reassembled from the text:
\[
\underbrace{\text{P1},\ \text{P2},\ \text{P3},\ \text{record individuation}}_{\text{consistency layer}} \;\longrightarrow\; \text{classification} \;\longrightarrow\; \text{residue};
\]
\[
+\ K \;\longrightarrow\; \text{Hilbert representation} \;\longrightarrow\; \text{Born statistics} \;\longrightarrow\; \text{projection, composition, evolution};
\]
\[
+\ \text{the realization bridge} \;\longrightarrow\; \text{physical unrealizability, and the irreversibility of observer-grade records.}
\]
The third line is the one the paper does not discharge, and it is where the formal result stops short of the physical one, (Section~\ref{sec:costs}).

\section{Setup and Postulates}\label{sec:postulates}

Fix a set of \emph{attributes}: questions with answer spaces (which slit was traversed; which detector fired). An \emph{information state}\footnote{A naming note, recorded once: the \emph{informational} in this classification is not to be confused with the Shannon-theoretic kind --- the states are logical objects, stated content, deductive closure, and blanks, with no entropy or channel notion in their specification, the quantitative layer entering only at the flagged kernel. It is to establish exactly this that the fuller development gives the states a proper name, \emph{Aethus}, plural \emph{Aethae} \cite{Aethus}.} $A$ is a deductively closed partial assignment of answers: $A(q) = v$ when attribute $q$ is \emph{stated} in $A$ with value $v$, and $A(q) = \ast$ when $q$ is \emph{blank}. Deductive closure means the state contains its own consequences. Nothing here is mental: an information state is carried by any record-bearing structure --- a collimator as readily as an experimenter --- and ``observer'' below means exactly such a structure. ``Observer-grade,'' used later for the record-bearing structures whose records the framework's physical claim concerns, names a subclass of these whose exact characterization is left open here; it is not defined as the class whose records cannot be kept from escaping, since that property is the claim about the class and not its definition, and the paper's conditional results concern specified record-bearing laboratories rather than a universal class. A \emph{proper child} $B \le A$ is an assignment extending $A$ by resolving any set of blanks along $A$'s own standing alternative structure ($B = A$ included), never by inserting content the structure does not offer, its cells being terms of $A$'s decomposition at coefficients not measure-zero relative to the rest, which at this paper's real weights is positivity. The relation is reflexive and transitive, and a state's proper children are all its descendants at every depth: any number of resolutions compacts into one, the blanks resolved along a chain being one set of blanks resolved at once, so there is no single-level proper child and every quantifier over $\le$ in this paper ranges over descendants. The canonical tree of states is ordered by this relation, the long paper's proper-child relation, a partial order; the results below use only that relation and the visible record cells of a configuration, and no meets or joins; a record partition carves a state into its cells with structural parent-conditional coefficients, $A \rightsquigarrow \sum_j \gamma_j (A \cap E_j)$, a declared decomposition and not an identity of the weighted algebra, written $\gamma_j$ to keep them apart from the kernel weights $W$, so record cells are components by construction and proper children where their coefficients are visible, with an edge from $A$ to every proper child of it, so that two consecutive resolutions $A \to A \cdot A_1 \to A \cdot A_1 \cdot A_2$ are the single resolution $A \to A \cdot (A_1 A_2)$, the product of the added content being itself one state; no privileged step-size exists for an edge, which is the no-privileged-decomposition principle of the companion's Attribute Set, and blanks are what a resolution resolves rather than what an edge counts. Every proper child is a \emph{refinement} of its parent, $B \sqsubseteq A$, refinement being the child relation $B = A \cdot C$, and proper childhood is the subrelation of refinement made of the visible terms of the parent's decompositions. In the companion's weighted setting, where the canonical statement of the pruning postulate lives, transitivity follows by composing decompositions, the relational weights multiplying along a chain, and the relation used here is its case with every cell visible \cite{Aethus}.

\begin{postulate}[Indexing]\label{p1}
Every determinacy attribution has the form ``$q$ takes value $v$ relative to $A$.'' There is no $A$-free ledger of determinate facts.
\end{postulate}

\begin{postulate}[Blanks]\label{p2}
Relative to any valid information state $A$: attribute $q$ is in superposition relative to $A$ if and only if $A(q) = \ast$. Blankness is superposition --- a complete informational entry in its own right, not a pointer to a hidden determinate value.
\end{postulate}

\begin{postulate}[Pruning]\label{p3}
An alternative survives only if it can be consistently refined; alternatives failing this are \emph{invalid} --- pruned, with no residual standing. Formally, validity $\mathcal{V}$ is the \emph{largest} predicate on states that is disjoint from the base contradictions and satisfies
\begin{equation}\label{eq:p3}
\mathcal{V}[A] \;\Longrightarrow\; \exists\, B \le A,\ \forall\, C \le B,\ \mathcal{V}[C],
\end{equation}
that is, the greatest fixed point of the safety operator, the complement of the doomed set of Eq.~\eqref{eq:mu} below. Read as a game: the observer refines, nature refines in reply, and $A$ is valid exactly when the observer has a refinement that survives every reply. Eq.~\eqref{eq:p3} alone is a necessary condition, satisfied by the empty predicate; the maximality clause is what makes it a definition.
\end{postulate}

Postulates \ref{p1} and \ref{p2} are stipulative; the two-quantifier alternation of Postulate \ref{p3} is not, and Section~\ref{sec:derivation} derives its form by elimination. Two remarks fix the reading. First, the pruning is driven by \emph{stated information} --- records, correlations, anything an information state carries --- and by nothing else: not size, not mass, not anyone's gaze. That austerity is what the classification theorem exploits. Second, Postulate \ref{p1} forbids free export of contradictions between information states: a clash internal to one state cannot, by itself, invalidate another, on pain of smuggling an $A$-free background ledger back in through the plumbing. The derivation below is precisely the resolution of the tension between this prohibition and the evident fact that \emph{future} contradictions do constrain \emph{present} states.

\emph{A remark on descent, for readers of the framework's other wings.} The relation $\le$ is proper childhood: $B \le A$ when $B$ extends $A$ by any set of blanks resolved along its standing alternatives, including none, so that $\le$ is reflexive and transitive, and a state's proper children are all its descendants at every depth. It is order-free and path-independent because extensions along standing alternatives commute and compose, resolution being extension by an attribute-set, not an attribute-sequence. Every quantifier over $\le$ below ranges over this relation. A tree is drawn by resolving one blank per edge, a drawing convention over which nothing quantifies; a chain of edges is one proper child, and the quantifiers see it as such.

\section{Deriving the Pruning Law}\label{sec:derivation}



A single-quantifier rule fails in both directions. An existential rule, invalid if \emph{some} refinement is invalid, dooms every state with any invalid descendant, which is every state that can be refined at all; a universal rule, invalid only if \emph{every} refinement is invalid, lets a state survive on the strength of one refinement that is itself unsafe. The least rule that avoids both is the two-quantifier safety form: a state is valid when it has a declared refinement all of whose declared descendants are valid, $\mathcal{V}(A) \Rightarrow \exists B \le A\ \forall C \le B\ \mathcal{V}(C)$, which is the third postulate. One possible bad descendant therefore excludes nothing; what excludes a state is that every refinement still has a contradiction-reaching descendant, and the quantifiers range over declared descent, the resolution of a state's own blanks along its standing structure, not over arbitrary insertion of content. Dually, the doomed set is the least fixed point
\begin{equation}\label{eq:mu}
\mathfrak{D} \;=\; \mu X.\big(\, \mathrm{Base} \,\cup\, \{\, A \;:\; \forall B \le A,\ \exists C \le B,\ C \in X \,\}\,\big),
\end{equation}
with $\mathrm{Base}$ the non-recursive modes of invalidity, and validity its complement under the base-generated reading.

\begin{proposition}[The fixed point closes at once]
\label{prop:closes}
Let $D_1 = \mathrm{Base} \cup \{A : \forall B \le A,\ \exists C \le B,\ C \in \mathrm{Base}\}$. Then $\mathfrak{D} = D_1$ for any branching whatever. Equivalently, a state is doomed iff every proper child of it has a base contradiction beneath it, and valid iff some proper child of it has a cone free of base contradictions:
\begin{equation}\label{eq:direct}
\mathcal{V}[A] \;\iff\; A \notin \mathrm{Base} \ \wedge\ \exists B \le A,\ \forall C \le B,\ C \notin \mathrm{Base}.
\end{equation}
\end{proposition}





The general equivalence, under the base-generation reading, is proved in \cite{Aethus}. Define finally the \emph{reduced form} $\mathcal{R}(A)$: the maximal safely-refinable core of $A$, the residue of alternatives surviving the pruning, on which a weight measure will later live. Postulate \ref{p3} is the statement that every valid state has one. Two remarks fence the derivation's scope. The first is exhaustiveness, since elimination over candidates is elimination only if the candidates exhaust the field. The field is the unweighted quantifier-prefix rules on the refinement relation, fixed as such by the division of labor between binary validity and graded weight, which excludes any rule keyed to how much weight the contradiction-reaching refinements carry. Within that field a two-quantifier rule is fixed by its first quantifier and its second, each ranging over refinements at every depth, read on the doomed side. At the first quantifier an existential is excluded by the inheritance direction: a rule that dooms a state through one refinement dooms that refinement's innocent siblings, whatever they hold deeper, so that quantifier is universal or the rule is wrong. At the second quantifier a universal never fires, since waiting is always a refinement that is not yet a contradiction, so that quantifier is existential or the rule is idle. That leaves $\forall\exists$-doom, whose De Morgan dual is the $\exists\forall$-validity form of Eq. This elimination covers the two-quantifier forms examined; idempotence of the $\forall\exists$ operator establishes its closure, not the equivalence of every deeper quantifier prefix to it, and the uniqueness claim is restricted accordingly.~\eqref{eq:p3}; the other alternation, $\exists\forall$-doom---some proper child all of whose proper children are doomed---is not a dual of anything displayed but a rule excluded at the first quantifier by the inheritance direction. Deeper nestings are not new candidates: by Proposition~\ref{prop:closes} the displayed alternation is already closed, and a deeper nesting adds nothing it does not already contain. That statement is restricted to the two-quantifier forms examined; it does not establish that every deeper prefix reduces to them. And the canonicity of this operator's least fixed point is not bespoke: base-generated doom with $\mu$-closure is the standard safety-property construction, Knaster--Tarski supplying existence, so the framework selects the canonical object of its type rather than an invented one. Uniqueness over all conceivable consistency rules is not claimed; elimination over the natural candidates, plus type-canonicity, is. The objection this invites should be stated in its strongest form, since the answer is the paper's actual claim rather than a retreat from it: \emph{you have selected a rule space in which your preferred rule is canonical}. Just so---and what is offered is that the rule space is stated explicitly, that the unique member within it is derived rather than chosen, and that no physical uniqueness outside it is claimed. A reader who proposes a rule outside the class is not refuted by anything here; they are proposing a different framework, and the comparison would be between frameworks rather than within one.

\section{Agreeing and Disagreeing Superposition: The Two Composition Regimes}\label{sec:regimes}
The postulates induce a structural dichotomy among co-present alternatives, and it will carry the entire classification. Let a blank attribute enumerate alternatives $a_1, \dots, a_n$, with $A_i$ the alternative-states. Two regimes, distinguished by what the composition \emph{includes}, under the framework's canonical names---\emph{agreeing superposition} and \emph{disagreeing superposition}, the composition companion's terms, with ``joint'' and ``exclusive'' their explanatory translations \cite{AethusUnion}:
\begin{equation}\label{eq:regimes}
\text{agreeing (joint):}\;\; \bigcup_i a_i \qquad\qquad \text{disagreeing (exclusive):}\;\; \bigcup_i \Big( a_i \cap \bigcap_{j \neq i} \neg a_j \Big),
\end{equation}
the lower-case letters being pins of one attribute, each term of the exclusive form the state whose interior holds one pin and excludes the rest, and the joint form the pins composed within one interior. Blankness itself is superposition, by $P_2$, independently of the regime a blank's alternatives occupy; so ``superposition'' is the wider term, and agreeing, disagreeing, and the constrained cases between are its regimes---a broader usage than the physics habit of equating superposition with phase-coherent interference, and the breadth is the point.
The union principle's own content is exactly this contrast: the agreeing element is the union with the interior negations \emph{omitted}, the exclusive composition the union with them \emph{included} --- a difference in what the element-set contains, not in the status of any flag. The record-side tracker of the same dichotomy is the \emph{interior} attribute $e_{ij}$ carrying the joint standing of alternatives $i$ and $j$: the exclusive composition is the one an observer's state renders with $e_{ij}$ declared false --- the negation stated as an entry --- the agreeing one with $e_{ij}$ undeclared, nothing said either way, and the two characterizations should not be conflated --- one is about composition, the other about records, apples and oranges even as they track one dichotomy. (A vocabulary rule, stated once and kept: the dichotomy is logical, declared-false against undeclared, and it is definable before any kernel exists. Before the kernel is introduced in Section~\ref{sec:axiom}, this paper uses only \emph{agreeing} and \emph{disagreeing} for the composition structure and does not use \emph{coherent} or \emph{decoherent} as its own vocabulary; those words enter only once the kernel does, as quantitative properties of the physical realization, and ``interaction'' talk for the agreeing regime is pedagogical shorthand used only where the physics has earned it---which interaction class runs on an agreeing pair is the dynamics' verdict, not the logic's, and mechanical interaction and quantum interference are not to be conflated through the word.)
The placement of the negations is the entire content: mutual exclusivity lives in the \emph{interior} of each pocket as stated non-interaction, never as an exterior wall between pockets. The inversion of burden is the deeply nonclassical content, simply stated: because exclusion is interior stated structure, \emph{neglecting} to state it does not leave composition undefined --- it exposes the alternatives to or-mediated \emph{jointness} as the automatic consequence --- the open interior being the open door, with matter's interference the downstream \emph{use} of what stands opened. And the mechanism's scope is fully general, which is why the agreeing regime is absent from macroscopic experience: the pruning postulate acts on any stated record whatever attribute it concerns --- path is the laboratory's exhibit, not the principle's extent --- so any agreeing superposition co-stated with a record it jointly invalidates is pruned to the disagreeing regime. The record does not fill the blank in; it eliminates the joint composition, and a blank can survive the pruning as a blank --- over alternatives that are now exclusive. Macroscopic environments record wholesale; every macroscopic composition arrives pre-pruned; and classical \emph{or}-exclusivity is thereby the special case that ubiquitous record-keeping manufactures --- not a rival composition law, but the joint law meeting saturated records. Exclusivity is thereby \emph{manufactured} --- an achievement of stated structure --- while interaction among co-present alternatives is the unpruned default.\footnote{The notation invites a classical misreading worth blocking explicitly, since the concept beneath it is perhaps the framework's most alien. The two regimes are not two configurations of one classical substance --- a list of alternatives with an interaction flag toggled on or off. The union principle beneath this display holds that composition itself is two-fold at the root: beside classical conjunction stands a \emph{primitive union} --- the jointness of \emph{or} --- realizing disjunction physically as joint occupation of the block, co-presence with openness as its native character, usable but, like \emph{and} itself, not further analyzable; the doctrine's canonical statement and development are the dedicated extraction's \cite{AethusUnion}, surveyed in the overview \cite{AethusBrief}. That substance is the \emph{default} composition of co-present alternatives; classical mutual exclusivity is not primitive but manufactured, the nested achievement in which interior negations forbid interaction while permitting simultaneous presence. The familiar physical \emph{or} is thereby, in the fuller taxonomy, entirely a subcategory of the disagreeing case --- and the agreeing case is categorically alien to classical metaphysics, a sister taxon rather than a variant of anything within it. Logical realism survives the alienness: non-contradiction is preserved because disjoint attributes taken in \emph{intersection} are exactly what the framework prunes as invalid. The blank $e_{ij}$ therefore marks no absent classical flag; it marks the primitive substance operating unimpeded, the stated $e_{ij}$ marking its manufactured suppression.}



\section{The Classification Theorem}\label{sec:theorem}

\noindent\emph{Read forward: this theorem is clause (iii) of the recovery theorem of Section~\ref{sec:recovery} proved before the kernel enters --- observables derived from records --- and it is stated here on the postulates alone so that the recovery can consume it.}

Fix the laboratory case: $n$ slits and an arbitrary configuration of detectors, each detector firing on traversal of any slit in its designated subset.
% >>> EDIT Q3 (Fable 5, Aug 29 2026): the setup indexed per Postulate 1, with the two delayed-choice families sorted
Everything below is indexed, per Postulate~\ref{p1}: the configuration is stated content relative to the state doing the classifying, and the partition and bins are that state's. Where detector-presence itself stands blank, the classification runs pocket-wise over the configuration pockets, one partition per pocket, with the pockets' own regime read off Section~\ref{sec:regimes}'s criterion like any other --- which sorts the delayed-choice literature into two families. Where the setting is chosen and recorded and only the classifying state has not learned it, the cross content is stated and the configuration pockets are \emph{exclusive}: one partition obtains, unread. Where the setting is carried by a quantum control left uncorrelated with any record \cite{ionicioiu2011proposal}, the cross content is undeclared and the pockets are \emph{agreeing}: the partition stands in the agreeing regime, the kernel carries cross-partition terms, and no fact of which partition obtains is available to be unread. The theorem is indifferent to which family it is handed.
% <<< END EDIT Q3
Throughout this section a \emph{detector} is a device whose marker states for distinct slits are orthogonal, so that its signature is deterministic; a partial or inefficient tag, whose marker states overlap---slit-dependent efficiencies $\eta_1 \neq \eta_2$ give overlap $c = \sqrt{\eta_1\eta_2} + \sqrt{(1-\eta_1)(1-\eta_2)}$---has no signature in this sense and is the kernel's business, handled in Section~\ref{sec:twotag} as a marker overlap. Assign each slit its \emph{detector signature} --- the tuple recording which detectors would fire if that slit alone were traversed. A signature of the presently listed detectors need not contain every future record, so the signature used here is stipulated to be \emph{complete for the declared configuration}, including every record of the slit traversal accessible within the configuration's declared refinements, and \emph{persistent}, every declared refinement retaining a descendant stating each such record. Completeness concerns what the signature includes and persistence what remains available throughout descent; a refinement that closes an avenue of access is what separates them, and neither implies the other. Slits with identical complete signatures are then permanently indistinguishable, no record of the configuration separating them. The signatures partition the slits into equivalence classes $E_1, \dots, E_k$. Two guards on the setup. What a record \emph{is}, since the partition it induces draws the boundary and the classification theorem rests on it more heavily than on anything else the paper imports: the record-individuation lemma is the single load-bearing result taken from the companion rather than proved here, and a reader auditing the dependency structure should mark it as such: it is not a physical primitive added to the framework: a record of an alternative is any degree of freedom whose state is correlated with that alternative and is stated in some refinement of some information state, which is the record-individuation lemma read as a definition---distinctions are individuated by the records that could support them, and there is no independent notion of a ``physically real distinction'' for the definition to be circular against, since the first postulate has already made determinacy two-place. The long paper's Fundamental Theorem of Aethic Reasoning \cite{Aethus} states the same individuation from the other direction, and it should be read as a statement about distinctions and not about values: where no record can ever be taken there is no distinction among the alternatives to be drawn relative to that index, so nothing selects among them and they are performed jointly---the agreeing regime---which supplies no value to any attribute and only fixes the grain at which the attribute space exists for that index. What the framework does not supply is the \emph{dynamics} that make some degrees of freedom correlate and stay correlated; that is decoherence theory's business, and every boundary statement below is exact conditional on the physically realized record structure, not derived from nothing. Signatures quantify over \emph{all} correlated degrees of freedom, preparation included: a collimated beam aimed at one slit is itself a record licensing the singleton, and without that totality the theorem would misclassify ordinary preparations. And classifying a proper subset of a class as invalid uses one further principle, \emph{record-individuation} --- alternatives are individuated by possible records, so a distinction no information could ever support is not a distinction the attribute space contains. The companion derives it from Postulate~\ref{p1} together with the record discipline, as its Record-Individuation Lemma \cite{Aethus}; here it is consumed as that lemma, imported rather than re-derived, and it is what bin (ii) below uses.

\begin{lemma}[Proper children in a configuration]\label{lem:children} Let $A$ be a state that itself states a detector configuration $\mathcal{D}$ with signature partition $\{E_1,\ldots,E_k\}$, so that the proper-child relation is fixed by $A$'s own content and not by an external parameter. The proper children of $A$ are exactly those along that partition: a blank of $A$ is resolvable to a value only when some record in $\mathcal{D}$ affirmatively renders that value distinguishable from its alternatives, so that ``which slit $= a_i$'' is a proper child of $A$ when and only when $\{a_i\}$ is a signature class. In particular, where no record separates the members of a class, the singleton resolutions inside that class are not proper children of any state: they are not resolutions the configuration makes available, and Postulate~\ref{p3}, which quantifies over proper children, never reaches them. The object Theorem~\ref{thm:powerset} classifies is the claim ``the interior attributes of $T$ stand blank, and the alternatives outside $T$ are excluded'': stated content that a state can hold, whose first conjunct conjoins with a record resolving an interior attribute into a direct contradiction (bin iii), and whose second conjunct is an exclusion no record supports when $T$ is a proper part of a class (bin ii).
\end{lemma}
\noindent\textit{Proof.} Refinement is the addition of stated content, and stated content is, by the record-individuation lemma just cited, individuated by possible records; so a resolution exists as a child exactly where a record in $\mathcal{D}$ can support it, which is the partition. That the classified object is stated content and not a blank is what distinguishes ``the particle is in coherent superposition over $T$'' from the mere absence of a which-slit entry: the former asserts, of the interior attributes, that they stand blank relative to every refinement, and that assertion is what a record can contradict. \hfill$\square$

\noindent The lemma is what keeps the paradigm valid state valid: in the no-detector configuration every singleton lies in bin (ii) below, and without the lemma every proper child of the classifying state would have a bin-(ii) proper child and Eq.~\eqref{eq:mu} would doom the state that ought to survive. With it, the singletons are not proper children, and the theorem reads as intended. The qualitative classification is proved for a declared record structure in which every candidate record cell is visible, its coefficient the carving's structural parent-conditional weight, not measure-zero relative to the rest; the weights the kernel later supplies enter only afterwards, and the general weight-filtered extension requires the corresponding visibility assumptions.

\noindent\textit{Structural visibility and kernel weights.} The two layers carry different quantities, and in this model the first is fixed independently of the second: every record cell receives a positive structural coefficient $\gamma_j$, whatever its kernel weight. They come apart at destructive interference. A cell whose paths' kernel contributions cancel at a detector, a dark port, has kernel weight $W = 0$ and remains a structurally visible proper child. The classification says which clusters are admissible and the kernel how often outcomes occur, so a zero-probability outcome is an admissible cell on which conditioning is undefined (the positivity paragraph after Theorem~\ref{thm:recovery}). Tying the structural coefficients to the kernel weights, so that a kernel-null cell would be no proper child, is a modeling choice this paper does not make. Structural visibility governs admissibility; probabilistic support governs occurrence and conditioning, and the two supports are not identified in this model, so a cell of zero kernel weight can still enter the third postulate's descendant quantifiers. This structural-visibility assignment is held conditionally, as a model of its own, until its correspondence with the physical weighting is established: the coefficient it supplies is the algebra companion's relative component weight, which is not by definition the probability of realization \cite{AethusTree}. What relation must hold between the two layers is an open compatibility question (Section~\ref{sec:costs}), and Section~\ref{sec:correspondence} takes the concrete step: one finite interferometer model specifying its carriers, declared decompositions, structural coefficients, kernel weights and base contradictions, with a proof of which proper-child relations and validity judgments the two readings share (Proposition~\ref{prop:readings}).

\begin{theorem}[Powerset Invalidation]\label{thm:powerset} Assume the record structure of the declared configuration is \emph{complete and persistent}: each slit's signature includes every record of its traversal accessible within the configuration's declared refinements, and every declared refinement retains a descendant stating each such record. Then every subset $T$ of the $n$ slits --- every candidate claim ``the particle traverses exactly these slits in coherent superposition'' --- falls into exactly one of three bins: (i) \emph{valid}: $T$ is empty or equals one full equivalence class; (ii) \emph{over-specified} (invalid by record-individuation): $T$ is a nonempty proper subset of a single class, drawing a distinction no information could ever support; (iii) \emph{consistency-pruned} (Postulate-\ref{p3} invalid): $T$ spans two or more classes, combining alternatives some record can distinguish, so that a future refinement can expose the maintained agreeing superposition as contradictory --- doomed in the fixed-point sense of Eq.~\eqref{eq:mu}. The bins are mutually exclusive and jointly exhaustive, with the counts explicit: $k$ valid nonempty cases, $\sum_j (2^{|E_j|} - 2)$ over-specified cases, and the remainder consistency-pruned --- the identity $1 + k + \sum_j (2^{|E_j|} - 2) + |\mathrm{bin\,(iii)}| = 2^{n}$ fixing the pruned count exactly.
\end{theorem}

\begin{proof}
For nonempty $T$, the signature partition gives an elementary trichotomy: $T$ is a proper subset of one class; $T$ equals one class; $T$ meets two or more classes. In the first case $T$ distinguishes slits that share every possible record, so by record-individuation no attribute carries the distinction and the claim over-specifies. In the second, every branch of the maintained superposition produces the identical total signature, so no refinement along accessible records can expose a contradiction, and the state has a clean future --- valid (empty $T$ trivially so; with no detectors, $k = 1$ and the sole valid nonempty case is the full set). In the third, some detector's record separates a slit of $T$ in one class from a slit in another, so the coherence claim conjoined with that record's stating is a base contradiction, and the record family is \emph{persistent} over the candidate's declared descent in the sense that every refinement has a further descendant stating a separating record, which is the theorem's persistence hypothesis, a condition on what remains available throughout descent and distinct from completeness, which concerns what the signature includes; by the conditional record-exclusion lemma (Lemma~\ref{lem:record-exclusion-q}) this is membership in the doomed set $\mathfrak{D}$ of Eq.~\eqref{eq:mu}. Exclusivity and exhaustiveness are the trichotomy's own, and the counts follow by enumeration.
\end{proof}

\noindent The theorem is the combinatorial classification once that record structure is given. The physical determination of the structure, which records a configuration makes accessible and which of them persist through its refinements, is the realization's task and not the theorem's.



Degrading a detector rather than removing it yields the first quantitative statement, still \emph{amplitude}-free: classically random partial tagging of the paths --- the tagger fires on a known fraction $q$ of runs, an exclusive-regime apparatus superposition over fired/not-fired --- produces fringe visibility $V$ and which-path distinguishability $D$ on the linear chord
\begin{equation}\label{eq:chord}
V + D = 1.
\end{equation}
The derivation is three lines. The run-ensemble is pruned into the tagged fraction $q$ --- paths distinguished, the multi-slit subsets consistency-pruned, no fringes, full which-path record --- and the untagged fraction $1-q$ --- signatures merged, the full class valid, full fringes, no record. Visibility and distinguishability are then the classical convex mixtures $V = (1-q)\cdot 1 + q\cdot 0 = 1-q$ and $D = q\cdot 1 + (1-q)\cdot 0 = q$, whence $V + D = 1$ identically. And the pairing of weight to element is not a bookkeeping choice: the tag's presence on a run is a stated attribute, so the regime dichotomy of Section~\ref{sec:regimes} places the tagged and untagged runs in the disagreeing regime, and $q$ rides the distinguished-paths element and $1-q$ the agreeing one because those are the two classes of that partition. The grade is not an input either: decided tags offer no amplitudes, so the surviving weights can only compose by the disagreeing regime's own convex law, as the overlapping polarizer tags of the two-tag configuration will put amplitudes back. Classical convex weights over exclusive alternatives are used --- the regime's own composition law --- together with the optical assumptions stated here and not elsewhere: an ideal untagged pattern at $V = 1$, a fully tagged pattern at $V = 0$, compatible intensity normalization and fringe geometry, and an accessible flag for the tagged subensemble, so that the chord is a conditional mixture calculation; no amplitudes and no kernel enter, and consistency and convex weights alone do not supply the endpoint patterns, as the depaired weighting below shows. The chord, moreover, is the efficient boundary of a region theorem:
\begin{theorem}[Triangle characterization]\label{thm:triangle}
Under arbitrary classical randomization over \emph{orthogonal-marker} detector configurations and record-access operations --- all-or-nothing tagging per run, with any probability, on any subset, subsequent destruction of record access permitted; an inefficient tag with overlapping marker states is classical in the everyday sense but lies off the triangle and belongs to Section~\ref{sec:twotag} --- the achievable $(V, D)$ pairs are the triangle $V + D \le 1$, $V, D \ge 0$, with the chord attained when no acquired record is wasted.
\end{theorem}
\begin{proof}
The vertices are realized: $(1,0)$ untagged; $(0,1)$ per-slit tagged; $(0,0)$ by tagging followed by destruction of record access --- the pruning has occurred and validity acts on elements whole, so fringes do not return, while accessible distinguishability is zero --- the generic macroscopic case, records access-destroyed while traces persist in principle; a coherent reversal of the marker interaction, a joint operation on system and marker, would instead restore the blank and return the run to $(1,0)$; discarding or destroying the marker does not, since for any trace-preserving operation $\mathcal{E}$ on the marker alone $\mathrm{Tr}_M[(I \otimes \mathcal{E})\rho_{SM}] = \rho_S$ and the unconditional pattern is unchanged, and the recovery of fringes by conditioning on an eraser outcome is a third case, the sorted subensembles of the preceding paragraph; which of these the trigger question of open construction five admits is its territory, flagged as such. Randomization convexifies, filling the triangle. For the bound: each run is tagged (contributing $(0, d_r)$, $d_r \le 1$) or untagged (contributing $(1, 0)$); with tagged fraction $q$ and mean accessible distinguishability $\bar d$, $V = 1 - q$ and $D = q\bar d \le q$, so $V + D \le 1$, equality iff no record is wasted.
\end{proof}

\noindent The split into $q$ and $1-q$ is not imposed by hand: the tag's presence on a run is itself a stated attribute, so the regime dichotomy of Section~\ref{sec:regimes} applied to that attribute places the tagged and untagged runs in the disagreeing regime, and the convex weights are the classical weights over that partition. No regime is selected anywhere in the derivation; the record selects it.

\noindent The lens between chord and circle, $1 - D \le V \le \sqrt{1 - D^{2}}$, is therefore the kernel's signature: under the stated optical assumptions, classical randomization over the specified tagging configurations yields the triangle, a calculation that does not require evaluating the path kernel, and the axiom buys the lens. The comparison is exact because $D$ is one quantity on both curves, the accessible distinguishability at the placement in force fixed in Section~\ref{sec:twotag}: the triangle's $(0,0)$ vertex is the record destroyed, and the circle's rim is the optimal placement, where $D$ attains $D^{\star}$---not two differently defined quantities set side by side. One further pre-kernel result then has its address here, since the display of Section~\ref{sec:axiom} lists it: Bayesian conditioning is the depaired-weighting case of the same machinery --- assign classical weights over disagreeing survivors, state new information, and pruning-plus-renormalization over the surviving residue is conditionalization, with the worked display in the companion \cite{Aethus}. The slogan the construction earns is \emph{one residue machine, two weightings}: the depaired weighting yields classical probability, the kernel of Section~\ref{sec:axiom} yields the quantum statistics. The circle awaits the axiom.

\section{The One Axiom}\label{sec:axiom}

Everything above ran on the postulates, enumeration and the record structure, without evaluating the path kernel: structural coefficients and, under the stated optical assumptions, classical mixture weights appeared, but no kernel weight did, and the chord and the triangle were amplitude-free. The quantitative layer is one axiom, the path kernel, and the discipline of this section is to state it and exactly what it supplies, so that the representation theorem's hypotheses can be read off rather than reassembled. The axiom assigns to each path its kernel,
\begin{equation}\label{eq:kernel}
W_x(T) \;=\; \Bigl|\sum_{p \in T,\; x(p) = x} \mathcal{K}(p)\Bigr|^{2}, \qquad W(T) \;=\; \sum_{x \in \Sigma} W_x(T) \;=\; \sum_{x}\ \sum_{p, q \in T:\ x(p) = x(q) = x} \mathrm{Re}\bigl[\mathcal{K}(p)\,\overline{\mathcal{K}(q)}\bigr],
\end{equation}
and nothing else: the weight of a cluster is the squared modulus of the summed kernel over its paths, which is the content the clauses below unpack.





\emph{The axiom's stated content.} Three clauses belong to the axiom and are not derived from it, and the recovery theorem below is graded against them. (K1) \emph{The weight law and its polarization}: a cluster's weight is Eq.~\eqref{eq:kernel}, equivalently the vector-valued assignment of Eq.~\eqref{eq:vector} with squared-norm weights;\footnote{One consequence of (K1). Because the kernel is installed as the pair sum over $T \times T$---the Born rule's expansion written out in full, cross-terms included---no amplitude is ever an object any information state holds: complex state-vector amplitudes are not primitive contents of any information state---the primitive quantitative object is the scalar pair-weight structure, of which the kernel's per-path value---the phase $e^{iS[p]/\hbar}$ in the scalar-action case---is an ingredient and not a state's content---and the state vector is needed only as the bookkeeping Lemma~\ref{lem:polarization} recovers from the weights at each configuration, and across configurations only as far as Claim~\ref{claim:global-polarization} reaches. Probability amplitudes are not posited and then interpreted; the weight law is the whole quantitative input, and amplitudes are what it factors into when one wants a Hilbert space.} Lemma~\ref{lem:polarization} shows the scalar weights suffice to reconstruct the kernel values at each configuration, the reconstruction across configurations being Claim~\ref{claim:global-polarization}, with three steps open---the moduli and relative phases together from the singleton and pair weights, which fix the full real Gram matrix of the kernel values---with the kernel's cross-slice composition structure available to tie orientations across configurations, so the one axiom is, said plainly, the path-integral state and its composition law. (K2) \emph{Unitarity}: the propagator $K(x',x) = \sum_{\text{routes}}\mathcal{K}$ is unitary on the physical state space. This is posited as part of the axiom's content: a real action and a normalized path measure do not by themselves make the finite propagator's columns orthonormal, and no derivation of unitarity from them is claimed; for the free internal kernel it follows from the resolution of the identity on the sphere. The composition law $\mathcal{K}(p_1 \cdot p_2) = \mathcal{K}(p_1)\mathcal{K}(p_2)$ alone is multiplicativity along concatenation---additivity of the action in the scalar case---and gives a semigroup; an absorbing (complex) potential satisfies it with a non-unitary propagator, so unitarity is posited here and not proved. With unit modulus no longer required of each path, it is (K2), and not the modulus, that separates an admissible internal kernel from an absorbing one. (K3) \emph{Calibration}: $W_x(T)/W(T)$ is the empirical frequency of the readout $x$ on runs prepared in the cluster $T$. The axiom is scoped to systems admitting such a compositional kernel---the scalar-action kernel over a classical action and, for polarization, the discrete SU(2) coherent-state kernel on the sphere, which supplies the free-propagation overlaps (Section~\ref{sec:polarized}); general internal dynamics and the continuum internal construction are among the open items, and the Dirac--von Neumann formalism is more general than the systems the axiom reaches.

The axiom's purchases are then theorems rather than inputs. The sequential-polarizer chain --- Malus's per-polarizer $\cos^2$/$\sin^2$ weights, interposed-polarizer revival included --- falls out as pruning residue once outcomes are fixed as determinate two-node overlaps, with the per-outcome decomposition and the polygonal geometric-phase structure of the standard sequential-measurement treatment arriving as consequences of the kernel rather than insertions. An interposed polarizer, in particular, does not manufacture a new state with freshly assigned amplitudes: it contributes a further determinate node, changing which augmented routes survive, while the same kernel assigns the overlaps along the surviving node-chain. Crossed polarizers at $0^\circ$ and $90^\circ$ admit no route, the overlap $\cos 90^\circ$ vanishing; with a $45^\circ$ polarizer between them, $0^\circ \to 45^\circ \to 90^\circ$ is a different admissible node-chain, with product amplitude $\cos 45^\circ \cdot \cos 45^\circ = \tfrac12$ and transmission $\tfrac14$. Nothing is restored to an earlier wave; the experiment has acquired a different boundary condition. And the complementarity budget splits exactly along the axiom line: the chord Eq.~\eqref{eq:chord} needed no kernel, while coherent partial tagging --- one agreeing-regime joint cluster --- reaches the circle, and the deductive step from the cluster to the identity is the marker expansion, the amplitude-grade analogue of the coupling lemma (writing $c$ throughout for the marker overlap): expanding one marker in the other's basis, $m_B = c\, m_{\parallel} + \sqrt{1-c^2}\, m_{\perp}$, the collinear component carries the interference and the orthogonal component the record, so the unconditional visibility is $V = c$ directly, and with the record-placement supremum evaluated as the geometry indicates and the companion derives \cite{Aethus}, $D^{\star} = \sqrt{1-c^2}$, so that, at every placement and with equality at the optimal one,
\begin{equation}\label{eq:circle}
V^2 + D^2 \le 1, \qquad V^2 + (D^{\star})^2 = 1 \ \text{for pure markers},
\end{equation}
the Pythagorean identity of that expansion, with equality for pure tagging at the optimal placement \cite{greenberger1988, englert1996},\footnote{The expansion, unpacked. $m_B = c\, m_{\parallel} + \sqrt{1-c^2}\, m_{\perp}$ is the rewriting of one marker in the other's basis and nothing more: its two legs carry weights $(1 + c^2)/2$, a tilted pair of paths that interferes, and $(1 - c^2)/2$, a single path that does not, and the total visibility is $c$.}

\section{A Hilbert Representation on the Finite Path-Kernel Domain}\label{sec:recovery}

% >>> ADDED (Fable 5.1, Sep 2 2026, author-directed re-centering): the recovery theorem ported from the
% companion into this paper's neutral vocabulary (information state; agreeing / disagreeing regime; the
% kernel), with the polarization lemma, the axiom-by-axiom table, and the placement against the
% reconstruction programs. This is the paper's principal result; the exhibits that follow are its clauses
% shown at work.
With the kernel installed, the formal package is closed---its named axiomatic core the three postulates and the three-clause kernel axiom, with record individuation, the composition primitive, and the product structure of composite setups as explicit background structures, all inventoried below, and the realization bridge standing outside the formal theorem---and the question is what it recovers. The answer is the standard formalism, axiom by axiom, and the recovery is a theorem. Fix the path space $X$ of a setup, (spacetime routes carrying whatever internal routes the setup's markers require), the compositional kernel $\mathcal{K}(p)$ of Eq.~\eqref{eq:kernel}, with $e^{iS[p]/\hbar}$ its scalar-action specialization, and a \emph{readout slice}: the set $\Sigma$ of configurations at which the setup is read, each $x \in \Sigma$ comprising the readout position together with every record-bearing degree of freedom the configuration carries---detector states, tags, the environment's registers---under the totality guard of Theorem~\ref{thm:powerset}. Every path $p$ ends at one configuration $x(p) \in \Sigma$. The Hilbert space is $\mathcal{H} = \ell^2(\Sigma)$ with orthonormal basis $\{|x\rangle : x \in \Sigma\}$ indexed by \emph{configurations}, not by paths, (the continuum case, with $|x\rangle$ delta-normalized and the sums replaced by the configuration and path measures, lies outside the theorem's rigor grade, as the classification's continuum extension does; the finite-slice statements are what is proved). Assign to every cluster $T \subseteq X$ the vector
\begin{equation}\label{eq:vector}
|T\rangle \;:=\; \sum_{p \in T} \mathcal{K}(p)\, |x(p)\rangle , \qquad \langle x|T\rangle \;=\; \sum_{p \in T,\; x(p)=x} \mathcal{K}(p) .
\end{equation}
A configuration is the finest grain at which the block records the readout, not the grain at which the display shows it: under the totality guard the absorption site is recorded in the detector's microscopic state even when the readout reports only a cell, so two endpoints within one display cell are record-separated and stand in the disagreeing regime. A readout \emph{cell} $C \subseteq \Sigma$ is therefore an exclusive union of configurations, and its probability on a cluster is the sum $\sum_{x \in C} W_x(T)$ of the endpoint weights, never the modulus of a summed amplitude---the two differ, and only the sum agrees with the fringes a pixel of finite width actually integrates. Within-cell coherent addition is correct when the detector's effect is a rank-one projector onto an extended mode, and that is a fact about the detector's physics that the guard records, not a consequence of Postulate~\ref{p2}. Paths ending at the same configuration add coherently; paths ending at different configurations are orthogonal. This is the structure the classification demands and the reason the basis must be the configurations: two paths with different readout configurations are separated by the readout record and stand in the disagreeing regime, while paths sharing a configuration carry no record separating them and stand in the agreeing one. A path basis would make every pair of paths orthogonal and delete the interference the theorem recovers.

\paragraph{The assumption package, and each theorem's signature.} Everything the theorem consumes is collected once, so that ``three postulates and one kernel'' is read as the informal count of named axioms it is and not as a count of primitive structure:
\[ \mathsf{Q} \;=\; \langle\, P_1, P_2, P_3,\ \mathsf{Rec},\ \mathsf{Comp},\ K_1, K_2, K_3,\ \mathsf{Bridge},\ \mathsf{Prod} \,\rangle, \]
with $\mathsf{Rec}$ record individuation, $\mathsf{Comp}$ the composition primitive of the union move, and $\mathsf{Bridge}$ the realization principle from formal invalidity to physical unrealizability; and, consumed by clause (vi) alone, $\mathsf{Prod}$, the product structure of a composite setup, $\Sigma_{12} = \Sigma_1 \times \Sigma_2$---a kinematic stipulation of the setup specification, of a piece with the path space and readout slice themselves, and no part of the kernel. Write $\mathsf{Q}_{\log} = \{P_1, P_2, P_3, \mathsf{Rec}, \mathsf{Comp}\}$, $\mathsf{Q}_{\mathrm{ker}} = \{K_1, K_2, K_3\}$, $\mathsf{Q}_{\mathrm{phys}} = \{\mathsf{Bridge}\}$, and $\mathsf{Restr}$, the stipulated record-compatible restriction, listed on its own because it is an update assumption and not a postulate. Then each result carries its signature: $\mathsf{Q}_{\log} \vdash$ the classification theorem; $\mathsf{Q}_{\log} + \mathsf{Q}_{\mathrm{ker}} \vdash$ the Hilbert representation and the Born statistics, with $\mathsf{Prod}$ added for composites; $\mathsf{Q}_{\log} + \mathsf{Q}_{\mathrm{ker}} + \mathsf{Restr} \vdash$ the L\"{u}ders representation of record-compatible restriction, $\Phi(T \cap E_j) = P_j \Phi(T)$, representing a supplied restriction by a projector without thereby selecting restriction as the physical update; and $\mathsf{Q}_{\log} + \mathsf{Q}_{\mathrm{ker}} + \mathsf{Q}_{\mathrm{phys}} \vdash$ the physical measurement interpretation. The rule $\mathcal{V}[A] \iff \exists B \le A\ \forall C \le B\ \mathcal{V}[C]$ is this paper's own construction, with $\le$ the proper-child relation, the admissible resolution of blanks that never adjoins contradictory content, so it does not inherit the algebraic companion's collapse, which concerns closed-content inclusion. That the companion's normal form establishes the same rule is a correspondence now specifiable but not yet proved \cite{AethusTree}: the companion's validity runs over a declared descent, the splittings each Aethus admits, and this paper's refinement is that descent as nature declares it: a blank that a record separates holds the carving the record carries and splits as an exclusive blank does, a filling asserted without a carving declaring no split, and a record-free agreeing superposition is irreducible, declaring no split until a record grants access, whereupon its splits are tagged, and so is it, which is the pruning of the agreeing cases this paper's applications turn on, so that the correspondence is intended over the whole declaration, exclusive and agreeing cases together; the embedding of this paper's information states into the companion's algebra is the declaration nature's Aethus carries, which this paper reads off record accessibility; what remains is the proof that the two validity predicates agree on that image. Until it is given, the rule stands on the construction given here.

\paragraph{Representation, faithfulness, completeness.} The map the theorem is about is $\Phi : T \mapsto |T\rangle = \sum_{p \in T} \mathcal{K}(p)\,|x(p)\rangle$, and three statements about it are logically distinct. \emph{Soundness}: every admissible cluster of nonzero weight, $W(T) > 0$, is assigned a ray, $T \mapsto [\Phi(T)]$. \emph{Faithfulness}: $[\Phi(T)] = [\Phi(T')] \Rightarrow T \sim T'$ for the appropriate framework equivalence. \emph{Completeness}: every ray of $\mathcal{H}$ is $[\Phi(T)]$ for some admissible $T$. The theorem proves the first. It does not prove the second, and the second is in general false as stated: distinct clusters can share a vector through coherent cancellation, so the exact object is the quotient $\mathcal{C}^{+}_{\mathrm{val}}/\!\sim_{\Phi}$ of the valid clusters of nonzero weight, $\mathcal{C}^{+}_{\mathrm{val}} = \{T \in \mathcal{C}_{\mathrm{val}} : W(T) > 0\}$---the empty cluster being valid with $\Phi(\varnothing) = 0$, and a nonempty cluster able to cancel to zero---, with $T \sim_{\Phi} T'$ iff $[\Phi(T)] = [\Phi(T')]$, on which $\bar{\Phi} : \mathcal{C}^{+}_{\mathrm{val}}/\!\sim_{\Phi} \hookrightarrow \mathbb{P}(\mathcal{H})$ is injective by construction; quantum-state equivalence is that quotient, represented faithfully by rays. The third has two statuses that must be kept apart. On a fixed finite path domain, with clusters as subsets and a fixed kernel, there are finitely many represented rays, so for Hilbert dimension at least two full ray completeness is impossible; completeness after enlarging the preparation domain is open, as clause (i) records. ``States are rays'' is therefore read as ``every valid agreeing cluster determines a ray,'' and nothing stronger.

\begin{theorem}[L\"{u}ders representation of record-compatible restriction, with the Dirac--von Neumann structure on the finite path-kernel domain]\label{thm:recovery} For a system admitting the compositional kernel of Eq.~\eqref{eq:kernel}, on the finite readout slices for which the theorem is rigorous, under the assignment \eqref{eq:vector} the axiom set yields the Dirac--von Neumann structure clause by clause, with every computed statistic agreeing; the continuum extension, general spin and internal dynamics together with the continuum internal construction, and the completeness of the state and observable representations on an enlarged preparation domain are outside the theorem's grade and are named among the open constructions, completeness on a fixed finite path domain being impossible for Hilbert dimension at least two. The finite free-polarization case the two-tag exhibit uses is inside it, its internal factor being the discrete SU(2) instance of the kernel. Each clause's inputs are those of Table~\ref{tab:kernel-boundary}: (iii) is a theorem of Postulates~\ref{p2}--\ref{p3} given the assignment, (v) represents the stipulated restriction $\mathsf{Restr}$, and the rest are read off the kernel axiom and the assignment.
\begin{enumerate}
\item[(i)] \textbf{Hilbert representation of cluster states.} A valid cluster in the agreeing regime is represented by the vector $|T\rangle$, and, when its weight is nonzero, $W(T) > 0$, by a ray. Its weight at a configuration $x$ is $W_x(T) = |\langle x|T\rangle|^2$, that is, $\bigl|\sum_{p \in T,\, x(p)=x}\mathcal{K}(p)\bigr|^2$, the coherent sum over the paths sharing that configuration---the kernel's display for the double slit at a screen point---and its total weight is the squared norm $W(T) = \sum_x W_x(T) = \langle T|T\rangle$; the state is the ray, normalized by $\sqrt{W(T)}$.
\item[(ii)] \textbf{Additivity under cluster union.} For disjoint clusters, $|T_1 \cup T_2\rangle = |T_1\rangle + |T_2\rangle$: composition in the agreeing regime is vector addition at kernel grade, and the cross-terms of $W(T_1 \cup T_2)$ are the interference the regime admits. This is the superposition principle at coefficients $(1,1)$; the representation's limits are collected after the proof. \item[(iii)] \textbf{Observables are self-adjoint operators, and they come from records.} A detector configuration's signature partition $\{E_1, \ldots, E_k\}$ (Theorem~\ref{thm:powerset}), is a partition of the slits by their deterministic detector signatures, and it is not by itself a partition of the readout configurations: two unmarked slits and two screen pixels give four routes, one slit-signature class $\{L, R\}$ and two endpoint classes $\{L \to 0, R \to 0\}$ and $\{L \to 1, R \to 1\}$, each slit contributing to both. Three partitions are in play, of the slits, of the routes, and of the endpoints, and the recovery uses the last: a record map $r : X \to J$ on complete augmented routes, whose fibers $r^{-1}(j)$ are the endpoint classes, carries the signature structure to the readout configurations, and it is those fibers that $E_j$ denotes below. Explicitly, with $x : X \to \Sigma$ the endpoint map, $\bar r : \Sigma \to J$ the record labeling of endpoints, $r = \bar r \circ x$, $\Sigma_j = \bar r^{-1}(j)$ and $E_j = x^{-1}(\Sigma_j)$, the projector is $P_j = \sum_{\sigma \in \Sigma_j} |\sigma\rangle\langle \sigma|$ and the central identity $\Phi(T \cap E_j) = P_j \Phi(T)$ follows; the factorization $r = \bar r \circ x$ is a requirement on the record, that two routes with the same endpoint never receive different labels. The partition used in the recovery is the one restricted to the prepared universe $T_0$: its classes are the pieces $F_j := T_0 \cap E_j$, the classification of Theorem~\ref{thm:powerset} is applied to $T_0$'s own partition and not to the ambient one, so that a piece is a whole class of the partition in force and not a proper subset of an ambient class, and preparation changes the classes in exactly this way. The transition itself is stipulated, not selected by $P_3$: $P_3$ is a validity predicate and specifies neither the successor index, nor the preservation of surviving kernel contributions, nor the restriction map; once record-compatible restriction is assumed, the identity represents it as L\"uders conditioning, and that is the theorem's claim, coherence at the preparation index being distinguished from coherence maintained through every accessible refinement. The cluster entering the calculation is the prepared cluster itself, $T = T_0$, valid at the pre-record index, before the record structure is asserted; asserting the record is a transition to a new index that changes the state's specification, and the post-record pieces $T_0 \cap E_j = F_j$ are valid clusters of that new state rather than refinements of a maintained superposition the classification tags invalid, refinements of an invalid state being invalid. The projector identity $\Phi(T \cap E_j) = P_j \Phi(T)$ holds for every cluster $T$; the validity claim is the classification's, and it is made for the prepared cluster, a proper subset of a piece $F_j$ being over-specified. The endpoint classes define orthogonal projectors $P_j$ onto $\mathrm{span}\{|\sigma\rangle : \sigma \in \Sigma_j\}$ resolving the identity, $\sum_j P_j = I$, so that any assignment of values $a_j$ to the classes is the self-adjoint operator $\sum_j a_j P_j$; a record at a later node, reached by propagation, is the same construction transported by (vii), $U^\dagger(\sum_j a_j P_j)U$, which is how observables not diagonal in the configuration basis arise. Formally the clause defines a map $\mathsf{Obs} : \mathcal{R} \to \mathrm{SA}(\mathcal{H})$, $\mathsf{Obs}(\{E_j, a_j\}) = \sum_j a_j P_j$, from record partitions with value assignments to self-adjoint operators, a soundness statement whose converse, surjectivity, is open. A candidate cluster $S \subseteq T_0$ is classified by its membership, which the vector represents: $S$ is valid and nonempty iff $S = F_j$ for some $j$, and then $|S\rangle$ lies in the range $P_j\mathcal{H}$; over-specified iff $S$ is a proper part of one $F_j$; consistency-pruned iff $S$ meets two or more of them. The vector of a consistency-pruned cluster has components in two or more ranges unless coherent cancellation removes them, as in $|0\rangle + |1\rangle - |1\rangle = |0\rangle$, which is why the bins are stated on membership rather than on support.
\item[(iv)] \textbf{The Born rule.} The probability of outcome $j$ on a cluster $T$ is $W(T \cap E_j)/W(T) = \lVert P_j|T\rangle\rVert^2/\langle T|T\rangle$, the weight of the record-restricted cluster. Three things are distinct here and are kept so: the per-path kernel value---the phase $\mathcal{K}(p) = e^{iS[p]/\hbar}$ in the scalar-action case---is primitive kernel structure; the weight law $W(T) = \lVert\Phi(T)\rVert^2$ is the kernel's calibration (K3); and $\langle x|T\rangle$ is a derived Hilbert-state representation of the path sums. So the ontological thesis is that state-vector amplitudes are not primitive contents borne by any Aethus, and it is not that complex amplitudes are derived from purely real weights: complex state-vector amplitudes are representationally reconstructed, and complex kernel values are primitive.
\item[(v)] \textbf{The projection postulate as the L\"{u}ders representation of record-compatible restriction.} Stating the record invalidates $T$ as an agreeing-regime object when it spans classes and leaves the exclusive-regime superposition of the pieces $T \cap E_j$ at the weights of (iv); conditioning on the outcome is pruning followed by renormalization and equals $P_j\rho P_j/\mathrm{Tr}(P_j\rho)$, the L\"uders rule---for an \emph{ideal} projective measurement, after any additional outcome-dependent dynamics $U_j$ within $P_j\mathcal{H}$ has been factored out, since a record with projector $P_j$ can coexist with such dynamics without changing which outcome was recorded, and the record partition alone selects the conditioning component and not every physical instrument. In formulas: the record partition generates the conditioning map $\mathcal{C}_j(\rho) = P_j \rho P_j / \mathrm{Tr}(P_j \rho)$, whereas a general instrument with the same outcome probabilities is $\mathcal{J}_j(\rho) = \Phi_j(P_j \rho P_j)$, linear, completely positive and unnormalized, with conditional state $\rho_j = \mathcal{J}_j(\rho)/\mathrm{Tr}(P_j \rho)$ and $\Phi_j$ any channel. Confining $\Phi_j$ to the recorded sector gives the sector-preserving instruments, a unitary or a reset among them, but the output need not stay there: $K_0 = |1\rangle\langle 0|$ and $K_1 = |0\rangle\langle 1|$ have $K_j^{\dagger} K_j = P_j$, the computational-basis statistics, with each output in the opposite sector; what the stipulated record-compatible restriction yields, represented in the cluster space, is $\mathcal{C}_j$, and the theorem is that and not $P_3 \Rightarrow \mathcal{I}_j$ for every instrument, nor $P_3 \Rightarrow \mathcal{C}_j$: $P_3$ supplies validity, the restriction is supplied with the record structure, and the result says which part of measurement is record-theoretic given that restriction and which remains dynamical. The rule selected is L\"uders's and not von Neumann's original: for a degenerate $P_j$, coherence \emph{within} the class survives, because within-class distinctions are unrecorded by bin (ii) of Theorem~\ref{thm:powerset}, whereas von Neumann's rule decoheres the degenerate subspace. The selection between the two rules is made by the stipulated restriction, not by the postulates, and once the restriction is fixed the selection is checkable. Mixed states are exclusive-regime superpositions, $\rho = \sum_i w_i |T_i\rangle\langle T_i|/\langle T_i|T_i\rangle$, convex because weights over exclusive alternatives add.
\item[(vi)] \textbf{Composite systems are tensor products.} The path space of a joint system is the product of the factors' path spaces and its readout slice the product of the factors' slices, so $\mathcal{H}_{12} = \ell^2(\Sigma_1 \times \Sigma_2) = \mathcal{H}_1 \otimes \mathcal{H}_2$.
\item[(vii)] \textbf{Evolution is unitary --- posited, as (K2).} The kernel's composition over intermediate nodes, $\mathcal{K}(p_1 \cdot p_2) = \mathcal{K}(p_1)\mathcal{K}(p_2)$ for concatenated routes, makes the map from a cluster's vector on one slice to its vector on a later slice, $\langle x'|T'\rangle = \sum_x K(x',x)\langle x|T\rangle$ with $K(x',x)$ the sum of $\mathcal{K}$ over the routes from $x$ to $x'$, the Feynman propagator \cite{feynman1948}; that it is unitary is clause (K2) of the axiom, not a consequence of the composition law, which by itself yields only a semigroup.
\item[(viii)] \textbf{No-signaling is an identity, for a fixed record-signature partition.} $\sum_j \lVert P_j|T\rangle\rVert^2 = \langle T|T\rangle$, and weights adding across record-separated branches is the statement that $\sum_j P_j\rho P_j$ has the statistics of $\rho$ for every observable commuting with the partition.
The identity is proved for a fixed record-signature partition with its fixed projector family and an already specified measurement structure; an extension in which the state can alter the kernel, the records or the admissible clusters, as in a geometry-dependent record map, is an extension, and no-signaling for such a model is a separate operational condition, $P(b \mid x, y) = P(b \mid y)$ for a spacelike-separated freely chosen $x$, to be established for that model and not inherited from this clause. \end{enumerate}
\end{theorem}

\noindent\emph{Positivity.} Three kinds of object are distinct: formal carriers, which the classification sorts whatever their weight; positive-weight clusters, $W(T) > 0$, which alone determine rays; and zero-probability outcomes, on which conditioning is undefined. A cluster whose kernel values cancel to a zero sum is a formal carrier the representation does not reach, and every clause that assigns a ray, normalizes or conditions carries the positivity hypothesis.
\noindent\textit{Proof.} (i) and (ii) are read off (K1) and the assignment: a vector defined as a sum of kernel terms over an orthonormal configuration basis has squared-norm weights and adds over disjoint unions, which is (K1) as installed and not a theorem of the postulates; (vi) is read off the assignment together with $\mathsf{Prod}$. In detail, (i): $\langle x|T\rangle$ is the coherent sum over the paths of $T$ ending at $x$, by the assignment and the orthonormality of the configuration basis, so $|\langle x|T\rangle|^2$ is the sum of the real pair-kernel contributions $\mathrm{Re}\,\mathcal{K}(p)\overline{\mathcal{K}(q)}$ over the pairs of paths sharing $x$, and $\langle T|T\rangle = \sum_x|\langle x|T\rangle|^2$ by Parseval. (ii): the sum over a disjoint union splits. (iii): the $P_j$ are orthogonal because the classes partition the configuration basis; a cluster's vector lies in $P_j\mathcal{H}$ if every path of the cluster ends in the $j$th class---not conversely, since coherent cancellation can remove a sector's contribution: path contributions $|0\rangle$, $|1\rangle$, $-|1\rangle$ sum to $|0\rangle$ although the cluster spans both classes--- and the three cases are the three ways a subset of paths can sit against the partition of their configurations, as in the proof of Theorem~\ref{thm:powerset}; the transported form is (vii) applied to a later slice's partition. (iv): $W(T \cap E_j) = \sum_{x \in E_j}|\langle x|T\rangle|^2 = \lVert P_j|T\rangle\rVert^2$ by (i). (v): $T$ spanning classes is invalid as an agreeing-regime object by (iii) and Postulate~\ref{p3}; what stands after the record is stated is the exclusive-regime superposition of the class-restricted pieces at the weights of (iv), renormalizing the surviving piece is conditioning, and $P_j|T\rangle\langle T|P_j/\lVert P_j|T\rangle\rVert^2$ is its Hilbert-space form; the mixture statement is the definition of the disagreeing regime, exclusions declared and weights adding. (vi): the joint readout slice is the product of the factors' slices, $\Sigma_{12} = \Sigma_1 \times \Sigma_2$, and $\ell^2(\Sigma_1 \times \Sigma_2) \cong \ell^2(\Sigma_1) \otimes \ell^2(\Sigma_2)$, which is the tensor-product structure; nothing about the joint kernel is used. In particular the kernel of a joint path need not factor: temporal multiplicativity, $\mathcal{K}(p_1 \cdot p_2) = \mathcal{K}(p_1)\mathcal{K}(p_2)$ for successive segments, says nothing about the kernel of a pair $(p_A, p_B)$ of subsystem paths, which factors, $\mathcal{K}_{AB}(p_A, p_B) = \mathcal{K}_A(p_A)\mathcal{K}_B(p_B)$, only for dynamically independent systems; an interacting joint kernel does not factor, and is what generates entanglement. (vii) is the composition law of the path integral. (viii): $\sum_j\lVert P_j|T\rangle\rVert^2 = \langle T|\sum_j P_j|T\rangle = \langle T|T\rangle$, and for $O$ commuting with every $P_j$, $\mathrm{Tr}(O\sum_j P_j\rho P_j) = \sum_j\mathrm{Tr}(P_jOP_j\rho) = \mathrm{Tr}(O\rho)$. Statistics agree throughout because every statistic the construction computes is a weight, and every weight is a squared norm by (i). \hfill$\square$

\noindent\textit{What the kernel adds to the probabilities.} Ontologically, the probabilities the kernel supplies are the framework's relative Aethic weights, realized through centric unfolding, which the long paper treats as one kind of resolution for quantum and classical attributes alike, its Randomness Universality Principle \cite{Aethus}. What the kernel adds is the numerical weighting for quantum attributes; the conditions under which the draw law's probabilities are preserved across finer resolutions are stated with that principle there. Centric unfolding prescribes the eventual outcome probabilities of a specified resolution, its weights in the validity quotient, while leaving the finite-step process underdetermined within its admissible family; for quantum attributes the kernel supplies those weights.

\noindent\textit{The representation claims, one by one.} Existence, uniqueness, soundness and completeness are separate, and each claim has its own. Each positive-weight cluster determines a unique ray. The representation need not be injective: distinct clusters may determine the same vector, and proportional nonzero vectors determine the same ray; and no single finite domain reaches every Hilbert state. Union is vector addition for disjoint clusters at coefficients $(1,1)$; closure under arbitrary complex combinations is not claimed. Record partitions give projection-valued observables soundly, provided the record map factors through the endpoints; whether every self-adjoint operator of interest arises as a transported record partition is open. Restriction is represented as L\"uders conditioning, the restriction being stipulated rather than selected by the validity rule, and conditioning is undefined on an outcome of zero weight. And amplitudes are recovered from weights locally, at each configuration up to a common phase and a common conjugation (Lemma~\ref{lem:polarization}), the reconstruction across configurations having open steps.

\paragraph{A non-triviality check: the depaired weighting.} Which clauses the kernel is responsible for is best displayed by running the same assignment with the depaired weighting
\[
 W_x^{\mathrm{dp}}(T) \;=\; \sum_{p \in T,\; x(p) = x}|\mathcal{K}(p)|^2 ,
\]
each path weighted by its own squared modulus and no pairing---a weight of one per path in the scalar-action case, and the squared overlaps where internal factors enter. The classification of Theorem~\ref{thm:powerset} survives unchanged, since it never consulted the kernel; the chord and the triangle do not, since they consult intensity profiles the classification does not supply. Under the depaired weighting two equal-modulus paths meeting a screen point give $W^{\mathrm{dp}}_x = |\mathcal{K}_L|^2 + |\mathcal{K}_R|^2 = 2$ with no phase-dependent fringe, so $V = 0$, and with no marker $D = 0$: $V + D = 0$. The chord is a conditional optical result. It needs an ideal untagged interference pattern at $V = 1$, a fully tagged pattern at $V = 0$, compatible intensity normalization and fringe geometry, and an accessible flag identifying the tagged subensemble; under those conditions the mixture calculation holds, and it is not a consequence of binary consistency and convex weights alone. Clause (iii) survives---records still induce partitions---and (v) survives in form, pruning-plus-renormalization becoming Bayes; (ii) and (iv) become classical: the union of clusters adds weights rather than amplitudes and the record-restricted weight is a classical marginal. Everything that separates the recovered formalism from classical probability is therefore the pairing in (K1) and the unitarity in (K2), and nothing else; the postulates supply the same structure under both weightings, which is what ``pre-kernel'' means.

\paragraph{Why the Born rule is not a second axiom.} A reader used to the standard order will look for the Born rule as a separate posit and not find it, and the reason is the point at which the two axiom sets differ in direction rather than in count. Nothing in the present set refers to a bare amplitude: superposition is the union of clusters, dynamics is the kernel's composition law on paths, records are partitions, mixtures are weighted lists of clusters, and the amplitude $\sum_{p \in T}\mathcal{K}(p)$ appears only inside the weight functional, never as an object an information state holds. The weight law is nonetheless enough to recover the amplitudes, because a quadratic form determines its polarization.
\begin{lemma}[Local reconstruction by polarization]\label{lem:polarization}
Let $W_x(T) = |\sum_{p \in T,\, x(p)=x}\mathcal{K}(p)|^2$ be the weight of a cluster at a readout configuration, for all clusters and all configurations on all readout slices. Then at each configuration the weights determine the kernel values ending there up to a common phase and a common complex conjugation, and where those values are collinear up to a common phase alone. At each configuration $x$ the singleton and pair weights give the real Gram matrix of the kernel values ending there,
\begin{equation}\label{eq:gram}
G_{pq} \;:=\; \mathrm{Re}\,\overline{\mathcal{K}(p)}\,\mathcal{K}(q), \qquad G_{pp} = W_x(\{p\}), \qquad G_{pq} = \tfrac12\big(W_x(\{p,q\}) - W_x(\{p\}) - W_x(\{q\})\big) \quad (p \neq q),
\end{equation}
which in the unit-modulus case reads $W_x(\{p,q\}) = 2 + 2\cos(\phi_p - \phi_q)$; and paths ending at different configurations are related covariantly, through the composition law (vii), by the weights at a later slice of clusters whose continuations recombine. The identity $2\,\mathrm{Re}\,\overline{\langle x|T_1\rangle}\langle x|T_2\rangle = W_x(T_1 \cup T_2) - W_x(T_1) - W_x(T_2)$, for disjoint $T_1, T_2$, extends the same to clusters.
\end{lemma}
\noindent\textit{Proof.} Read each kernel value as a vector $v_p = (\mathrm{Re}\,\mathcal{K}(p), \mathrm{Im}\,\mathcal{K}(p)) \in \mathbb{R}^2$. Then $G_{pq} = \langle v_p, v_q\rangle$, so \eqref{eq:gram} is the real Gram matrix of the $v_p$ at $x$, and a Gram matrix determines its vectors up to an orthogonal transformation of $\mathbb{R}^2$. The residual group $O(2)$ is exactly the stated freedom at one configuration: its rotations are multiplication by a common phase, its reflections a common complex conjugation composed with a phase, and where the vectors are collinear a reflection is itself a rotation, so no conjugation bit remains.  No weight of three or more paths adds information, since $W_x(\{p,q,r\}) = G_{pp} + G_{qq} + G_{rr} + 2(G_{pq} + G_{pr} + G_{qr})$ is fixed by the singletons and pairs; and no genericity is needed at a single configuration, a zero-modulus path being the zero vector and a collinear family merely shrinking the residual freedom.  The identity is the expansion of $|a_1 + a_2|^2$ with $a_i = \langle x|T_i\rangle$. \hfill$\square$


\begin{claim}[Global reconstruction; three steps unproved]\label{claim:global-polarization}
Across configurations, the weights determine $\mathcal{K}$ up to the basis-phase gauge of the configurations and one common complex conjugation per connected component of the \emph{sign-transmitting graph}. Its vertices are the configurations whose kernel vectors span the plane; two are joined when some later slice's recombination has a cross term between clusters through them, with nonzero propagator entries, whose value changes when the amplitudes on one side alone are conjugated. A configuration whose vectors are collinear carries no conjugation freedom of its own, conjugation there being a rotation, and a single conjugation remains when the graph is connected. Three steps remain proof obligations: (G1) that the edge condition is invariant under the phase gauges; (G2) that each edge imposes exactly the parity constraint equating its endpoints' bits; and (G3) that no constraint links distinct components.
\end{claim}

\noindent\textit{Argument in outline, not a proof.} Paths ending at different configurations have no pair weight at either, so each configuration's values are fixed only up to its own $O(2)$ frame, and the two parts of that frame behave differently. Its rotation is gauge. Rephasing a configuration, $|x\rangle \mapsto e^{i\chi_x}|x\rangle$, sends $\langle x|T\rangle \mapsto e^{-i\chi_x}\langle x|T\rangle$ and $K(x'',x) \mapsto e^{-i\chi_{x''}} K(x'',x)\, e^{i\chi_x}$, so every concatenated path, and every weight at every slice, is unchanged, the global phase being the constant case; recombination at a later slice therefore relates the configurations covariantly, $K(x'',x)\langle x|T\rangle$ being invariant, and does not fix their phases. Its reflection is not gauge: conjugating the values at one configuration alone changes the concatenated values, and with them the weights wherever continuations recombine. A recombination cross term between clusters through $x$ and $y$ is the real part of a product of amplitudes carried from $x$ and from $y$; if conjugating one side alone changes it, the weights fix the two bits equal, and if not, the carried relative phase being $0$ or $\pi$ or a propagator entry vanishing, it fixes nothing, which is the degenerate case the definition of the graph excludes. Equality propagates along edges, so each component carries one bit, and distinct components are constrained by no weight. Zero-modulus paths are zero vectors and enter no edge. The outline takes (G1)--(G3) for granted, and every use of the claim below carries that grade: the local lemma is available wherever it is invoked, and uniqueness across configurations is not established until (G1)--(G3) are.

\noindent At a single configuration the residual freedoms are exactly the ones a quantum representation carries there, the $O(2)$ ambiguity of a common phase and, where it applies, a conjugation, which is why the lemma claims no more. Across configurations the reduction of the residual freedom to the componentwise one described next is conjectural, resting on the obligations (G1)--(G3) of Claim~\ref{claim:global-polarization}: the local freedom is established, its global reduction is not. Two directions of construction are distinct. A supplied complex kernel gives the represented state directly; the scalar cluster weights recover its local Gram geometry; and reconstruction of the state from the weights alone across configurations is the open part, on which the recovery theorem does not wait, the first direction sufficing for it. The basis-phase gauge is the ordinary phase convention of the configuration basis, the global phase its constant case; the common conjugation is an antiunitary orientation ambiguity, weights being unable to fix the orientation of the complex structure. In the scalar-action realization the conjugation is the choice between $e^{iS/\hbar}$ and its conjugate, fixed by the sign convention on the action and so, there, by the time orientation of the paths; for a general complex kernel with no action behind it, identifying the conjugation with time reversal would need more physical structure than bare conjugation supplies, and is not claimed. So the state vector of (i) is a \emph{derived} bookkeeping object, and where the standard presentation posits a state space and then a Born map, the present set posits one weight law and reconstructs the Hilbert representation on the represented state image, surjectivity of that image onto the projective space being excluded on any fixed finite path space and its completeness for an enlarged preparation domain being the further problem the representation paragraph states. This is the third of the three relocations named in the introduction, and it is the one with the longest reach: probability is the root and the amplitude is derived from it, which reverses the order the formalism has carried since 1926 and places the framework's quantitative primitive in the same position the companion's anthropic development gives its own---a weight, and never a census or a squared number beneath it \cite{Aethus}. The standard pair is not independent either --- Gleason's theorem derives the Born rule from the state space and non-contextual probability assignments --- so the difference is one of direction: there, states posited and probabilities derived; here, weights posited and the representation of the states reconstructed, and only the second direction is one on which the amplitude is nobody's content.

\paragraph{The two axiom sets, side by side.}
\begin{center}
\footnotesize
\begin{tabular}{@{}p{0.27\linewidth}p{0.42\linewidth}p{0.26\linewidth}@{}}
\hline
\textbf{Dirac--von Neumann axiom} & \textbf{In the present set} & \textbf{Standing} \\
\hline
states are rays in a complex Hilbert space & valid agreeing-regime clusters, $|T\rangle \in \ell^2(\Sigma)$ over the readout configurations, normalized by $\sqrt{W}$ & kernel content (K1) \\
the superposition principle & the union of clusters, $|T_1 \cup T_2\rangle = |T_1\rangle + |T_2\rangle$ & kernel content (K1) \\
observables are self-adjoint operators & signature partitions of records, transported by propagation & derived (Postulates~\ref{p2}--\ref{p3}, Theorem~\ref{thm:powerset}) \\
the Born rule & the weight of the record-restricted cluster & \emph{absorbed}: (K1) with the calibration (K3); amplitudes at each configuration by Lemma~\ref{lem:polarization}, across configurations by Claim~\ref{claim:global-polarization} \\
the projection postulate & representation of record-compatible restriction as pruning then renormalization, the transition to the record index stipulated & kernel side: needs the Hilbert assignment and the kernel, with $P_3$ supplying validity only; L\"uders's form selected \\
composite systems are tensor products & product of readout slices, no kernel factorization used & the representation of (K1) with $\mathsf{Prod}$ \\
unitary evolution & the kernel's composition law, unitarity posited & posited (K2) \\
\hline
\textit{(no counterpart)} & the index: every state relative to an information state & posited (Postulate~\ref{p1}) \\
\textit{(no counterpart)} & the two regimes, with the classical limit as the exclusive one & posited (Postulate~\ref{p2}), the limit derived \\
\textit{(no counterpart)} & measurement's origin: which partitions the records induce & derived \\
\hline
\end{tabular}
\end{center}
The count, read off the table, is this: two of the seven standard axioms---the observables axiom---is a theorem of the postulates once the kernel is installed, and one---the projection postulate---is the L\"{u}ders representation of a stipulated record-compatible restriction, relocated rather than eliminated; three---the ray representation on the represented image, the superposition principle, and unitary evolution---are the kernel axiom's own stated content, (K1)--(K2), and are not derived from anything; composite systems follow from that representation and the product slice of a composite setup, $\mathsf{Prod}$, a kinematic stipulation that uses nothing of the joint kernel; and the seventh, the Born rule, is absorbed under (K3). So the set is not much shorter by count---four axioms against seven---but shorter in kind: it has no separate measurement axiom, no projection axiom and no axiom introducing observables, the observables arriving as consequences of the record structure and the pruning law and the projection rule as the representation of a stipulated record-compatible restriction, an assumption relocated rather than removed; and it carries three things the standard set has no slot for.

\paragraph{Placement against the reconstruction programs, and what the set does not do.} The kernel does not follow from the three postulates, and it should not be expected to. The postulates' substance is logical --- indexing, blankness, pruning, the jointness of \emph{or} --- and their attribute side carries a binary flag, exclusions declared or undeclared, while the kernel carries continuous phase; no pruning of a binary structure produces a continuous one. In the vocabulary of the reconstruction programs, the present set without the kernel is a test space in the sense of Foulis and Randall --- attributes as tests, alternatives as outcomes, same-attribute exclusivity as orthogonality --- with one addition those spaces lack, the agreeing-regime union; and against the reconstruction axioms it sorts cleanly. Convexity and the existence of measurements are theorems of the postulates alone, and pure conditioning is the representation of the stipulated restriction, the existence-of-measurements result being the one respect in which this base is stronger than the reconstructions' starting point, which takes measurement as primitive; no-signaling is not among them, being an identity of the kernel's weight law and so on the kernel side of Table~\ref{tab:kernel-boundary}, where the table places it. Continuous reversibility, local tomography, and the quadratic form are the kernel's content and nothing else's. The kernel's substance is accordingly kinematic where the postulates' is logical --- it reports a geometric arena that the agreeing regime toggles visibility into --- and the two sectors meet at exactly one door, the agreeing regime; deriving the kernel from the postulates would be a category mistake of the kind that deriving a metric from a logic would be, and its derivation is assigned to a geometric program on that ground (Section~\ref{sec:costs}). One consequence for the order of generality should be stated, since the quantum wing tempts one to invert it: probability --- weight over exclusive alternatives, convex --- is the general structure, and probability amplitudes are what that structure looks like when the agreeing regime's clusters are the things weighted, a special case tailored to the arena's behavior rather than a widening of probability; the polarization lemma makes the subordination exact.
% <<< END ADDED

\section{The Representation, Exhibited}\label{sec:exhibits}

% >>> ADDED (Fable 5.1, Sep 2 2026): the three constructions below are the recovery theorem's clauses shown
% at work, and each subsection now opens by naming the clause it exhibits.
Three constructions show the theorem at work, and each is introduced by the clause of Theorem~\ref{thm:recovery} it exhibits. The polarized slit exhibits (v): projection as the residue pruning leaves. The two-tag configuration exhibits (ii), (iv), and (viii) at once: both composition laws in one experiment, the Born weights across the ports, no-signaling as an identity --- and the section after it shows why the recovered objects have the grain of the formalism, by three standing theorems. The Bell leg (Section~\ref{sec:bell}), exhibits what the recovery's realism amounts to against the theorem that prices every other realism.

\subsection{The Polarized-Slit Construction: Projection as Pruning Residue}\label{sec:polarized}

\noindent\emph{Exhibiting clause (v) of Theorem~\ref{thm:recovery}: the projection postulate as the residue pruning leaves.}

The two-tag configuration of Section~\ref{sec:twotag} is best entered through its one-polarizer ancestor, built from the geometry upward. Fix each polarizer's transmission outcome as a determinate node on the polarization sphere (axis angles doubled, orthogonal polarizations antipodal): each path becomes an augmented path --- a spatial route carrying a polarization route --- and the amplitudes come with the geometry, the transmission amplitude at incident polarization $\alpha$ against axis $\theta_j$ being the two-node overlap $\cos(\alpha-\theta_j)$ and the absorption amplitude its antipodal counterpart $\sin(\alpha-\theta_j)$ --- the \emph{weights}, per Eq.~\eqref{eq:kernel}, their kernel squares.

\paragraph{Where the overlaps come from.} The analyzer coefficients are not additional measurement postulates, and they are not applied to the kernel after the port outcome: they are finite propagator overlaps of the augmented kernel. On the augmented configuration space the kernel of a path factorizes into a spacetime factor and an internal factor, and the internal factor is built from one per-step kernel on the polarization sphere,
\begin{equation}\label{eq:augkernel}
\begin{gathered}
\mathcal{K}_{\mathrm{aug}}[p] \;=\; \mathcal{K}_{\mathrm{st}}[p_{\mathrm{st}}]\; \mathcal{K}_{\mathrm{int}}[p_{\mathrm{int}}], \qquad \mathcal{K}_{\mathrm{int}}[n_0, \ldots, n_N] \;=\; \prod_{k=0}^{N-1} \kappa(n_{k+1}, n_k), \\[2pt]
\kappa(n', n) \;:=\; \langle n' \vert n \rangle \;=\; \cos(\Theta/2)\, e^{\,i\Phi_{\mathrm{P}}},
\end{gathered}
\end{equation}
$\kappa$ being the SU(2) qubit coherent-state kernel (the spin-$\tfrac12$ representation) for polarization---the polarization sphere and the qubit Bloch sphere sharing their mathematics, with no claim that the photon carries spin $\tfrac12$---$\Theta$ the angle between the two nodes, and $\Phi_{\mathrm{P}}$ the Pancharatnam phase in the chosen coherent-state gauge \cite{pancharatnam1956generalized}, gauge-dependent for a single open overlap, its gauge-invariant content appearing only around a closed loop, as the Berry geometric phase \cite{berry1984quantal}. The path factor is multiplicative under concatenation, as (K1) requires. The endpoint overlap is not---$\langle n_3 \vert n_1 \rangle \neq \langle n_3 \vert n_2 \rangle\langle n_2 \vert n_1 \rangle$ in general---and is instead the endpoint propagator, recovered by integrating the path factor over the intermediate nodes,
\begin{equation}\label{eq:endpoint}
\int \prod_{k=1}^{N-1} d\mu(n_k)\; \mathcal{K}_{\mathrm{int}}[n_0, \ldots, n_N] \;=\; \langle n_N \vert n_0 \rangle, \qquad d\mu(n) = \tfrac{2s+1}{4\pi}\, d\Omega,
\end{equation}
which is what inserting the resolution of the identity on the sphere does, at every slicing, in the exact time-sliced SU(2) coherent-state construction \cite{Kochetov1995}; the per-step values and the endpoint overlap are one kernel under one composition law. The two-tag construction uses the direct tag-to-analyzer leg, $N = 1$, where path factor and endpoint propagator coincide: $\kappa(\theta, \pm\alpha) = \langle \theta \vert \pm\alpha \rangle = c_{1,2}$. Physical angles are doubled on the sphere, so $\cos(\Theta/2) = \cos(\alpha - \theta_j)$. The coherent transmission amplitude is $\kappa$ itself, sign and phase included---which is what separates a port's fringes from its anti-fringes---its modulus is $\cos(\Theta/2)$, and the Malus transmission probability is $|\kappa|^2$. The formal continuum action of the time-sliced construction contains the Berry term, $\hbar s \int (\cos\vartheta - 1)\, d\varphi$ at $s = \tfrac12$.

Three things are not claimed, so that the scope is exact. The overlap is not derived from the scalar spacetime action alone: the internal geometry---the coherent states, their overlaps, and the resolution of the identity---is supplied as kernel structure, as the spacetime phase is. The per-step factors $\kappa$ are not unit-modulus pure phases, since $\vert\langle n_{k+1} \vert n_k \rangle\vert = \cos(\delta\Theta/2) < 1$ at finite slicing; the pure-phase form $e^{i(S_{\mathrm{st}} + S_{\mathrm{B}})/\hbar}$ holds only as the formal continuum limit. The recovery does not need that form: (K1) is stated for a general compositional kernel and Lemma~\ref{lem:polarization} recovers moduli as well as phases, so augmented paths fall under both as they stand; the regularized continuum construction, with its known boundary and correction terms \cite{Kochetov1995}, remains open as a separate matter. And the construction covers free propagation between optical elements, which is what these exhibits use; general internal dynamics under a nonzero Hamiltonian, and the continuum internal construction, remain open. What is established is what the two-tag exhibit needs, and a little more: one augmented Nexus supplies $c_1$ and $c_2$, the analyzer does not reweight afterward, and the construction falls under the kernel axiom and the recovery theorem as stated rather than as an exception to them. The construction is then pure Section-\ref{sec:postulates} machinery: a Boolean transmission-or-absorption attribute per polarizer encounter, the joint alternatives enumerated in superposition with the overlap amplitudes applied consistently along each path, the kernel-inconsistent joint combinations pruned, and the residue delivering exactly the per-polarizer Malus factors that the projector formalism books by hand --- projection as the reading of a pruning residue at node grade. Classically random partial tagging of the paths inside this construction is what produced the chord $V + D = 1$ of Eq.~\eqref{eq:chord}; coherent tagging is what reaches the circle of Eq.~\eqref{eq:circle}. One feature of the construction, stated: determinate nodes are never posited --- they form by the pruning of all paths not passing through them, and the two-tag's surviving pair below is that same mechanism's or-generalization, pruning all paths not passing through one node \emph{or} the other. Three clauses complete the statement. First, \emph{which} positions stand determinate is fixed by the system state itself: the measurement positions --- slit apertures and screen in spacetime, analyzer settings on the Bloch sphere --- are part of the state's determinate content, so the node-set is predetermined per state exactly as the outcome structure is. Second, a node is itself or-shaped: it names a \emph{subset} of spacetime (an aperture, a detector face) or of the sphere, and the surviving condition is passage through \emph{some} point of that subset, not through a designated one --- the node is a disjunction the pruning leaves standing, not a point the theorist marks --- and, being residue rather than posit, it inherits whatever measurement-error fuzziness attends its generation: the subset's diameter is a physical dial, with the whole-space limit blanking the node entirely (Section~\ref{sec:beyond}). Third, the nodes are the boundary conditions of the whole construction, imposed by pruning rather than by decree: between them, the agreeing superpositions of paths span the intermediate spacetime unconstrained, and the surviving path-bundles are exactly those threading every node --- determinacy at the nodes, superposition between them, with the pruning holding both ends of that division at once.

\paragraph{The textbook reading, translated.} The construction reverses the usual explanatory order without changing any statistic. In the projector formalism a recorded outcome is represented by the state update $|\psi\rangle \mapsto P_j|\psi\rangle / \lVert P_j|\psi\rangle \rVert$, after which the updated vector supplies the probabilities of the next measurement. Here the same sequential statistics are read in the opposite direction: the record determines a partition, pruning retains the augmented routes compatible with the new determinate node, and the unchanged kernel supplies their transition weights. The kernel coefficients of the survivors are not edited, and normalization only rescales their common conditional weight; the projected vector is the Hilbert-space representation of the surviving residue, $|T \cap E_j\rangle = P_j|T\rangle$ with its normalization. At representation grade the two are the same object, so no empirically different collapse law is proposed: the disagreement is over which side is primitive, the standard formalism taking projection as the update rule and the present one taking the record-restricted residue as primitive and recovering projection as its representation. Two consequences are worth stating in a textbook reader's terms. Pruning changes membership, not amplitudes: a record removes incompatible routes and never reaches inside a surviving route to alter its kernel coefficient, the kernel changing weights only through physical path structure. And a record settles exactly the distinction it records and no further one: an earlier determinate outcome remains part of the history without becoming a hidden, persistent valuation of later incompatible attributes---which is why, in the two-tag configuration, a port can stand determinate while the slit stays blank. The three textbook words translate accordingly: \emph{collapse} is pruning plus conditioning; \emph{superposition} is the agreeing regime, the alternatives the records leave unresolved, composing at amplitude grade; and a \emph{measurement basis} is a record partition, the determinate-node structure a detector configuration induces---the dictionary the comparison table of Section~\ref{sec:recovery} lays out axiom by axiom.

The two-tag configuration is that node-grade construction run where it works hardest.

\subsection{The Two-Tag Configuration: Both Regimes in One Experiment}\label{sec:twotag}

\noindent\emph{Exhibiting clauses (ii), (iv), and (viii) of Theorem~\ref{thm:recovery}: both composition laws in one experiment, the Born weights across the ports, and no-signaling as an identity.}

\noindent\textit{The chain, step by step.} The experiment runs the whole construction; given the declared record partition, its admissibility classification does not require evaluating the kernel weights. (1) \emph{Record configuration}: two slits whose routes carry polarization markers at $\pm\alpha$, with overlap $c = \cos 2\alpha$, and an analyzer at angle $\theta$ whose pass-or-block outcome is a record. (2) \emph{Signature partition}: for the declared analyzer configuration, no additional persistent record perfectly separates the two routes within each port, so the slits form one signature class; the port outcome partitions the four augmented routes into the two port classes of Eq.~\eqref{eq:ledger}, and their unequal kernel amplitudes retain the partial distinguishability. (3) \emph{Admissibility classification}: each port class is valid, a cluster mixing the ports is consistency-pruned, and within each port the slit pair stands in the agreeing regime. (4) \emph{Kernel}: the surviving routes carry $\mathcal{K}(p_L^{+}) = c_1 e^{iS_L/\hbar}$ and $\mathcal{K}(p_R^{+}) = c_2 e^{iS_R/\hbar}$, with $d_1, d_2$ in the other port. (5) \emph{Represented state}: within each port the pair is the summed amplitude of Eq.~\eqref{eq:portamp}. (6) \emph{Detection probabilities}: the port weights $p_\theta$ of Eq.~\eqref{eq:ports}, the in-port fringes and visibility of Eq.~\eqref{eq:twotag}, and the $\theta$-independent unconditional pattern of Eq.~\eqref{eq:nosig}. (7) \emph{Conditional update}: conditioning on a port restricts the cluster, $T \to T \cap E_{+}$, the stipulated restriction represented as L\"uders conditioning, and renormalization by one common factor leaves the ratio $c_1 : c_2$ unchanged. What changes with the records is the partition and the classification of steps (2)--(3); what requires the quantitative axiom is steps (4)--(7), the marker overlap $c$ of step (1) already belonging to it.



The endpoints come first. With no which-path record available, $D = 0$ and $V = 1$: the two paths remain one agreeing pair. With a perfect record, $D = 1$ and $V = 0$: the paths are fully separated into the disagreeing regime. Between them there are two qualitatively different ways to be intermediate, and they have different shapes. A \emph{classical} partial tag---some runs recorded, some not---lies on the chord $V + D = 1$ of Eq.~\eqref{eq:chord}, which Section~\ref{sec:regimes} derived in three lines as a convex split: a recorded component at $(D, V) = (1, 0)$ and an unrecorded one at $(0, 1)$, composed by weights $q$ and $1 - q$, so that $D = q$ and $V = 1 - q$. The straight line is therefore not an additional complementarity postulate; it is the geometry of mixing an interference-bearing agreeing sector with a record-bearing disagreeing sector, and it needs the convex weight structure and nothing from the kernel. A \emph{coherent} partial tag is intermediate in a different way: every run carries both path alternatives together, and what varies is how alike their marker states are. Writing $c$ for the marker overlap, complete overlap is no record ($D = 0$, $V = 1$), orthogonal markers are a perfect record ($D = 1$, $V = 0$), and partial overlap gives intermediate values that do \emph{not} lie on the chord, because the alternatives were never split into record-separated components---they remain one agreeing pair whose markers merely fail to coincide. For pure markers, with $V$ the \emph{unconditional} visibility and $D$ evaluated at the optimal placement, where it equals its supremum $D^{\star}$,
\[ V = c, \qquad D^{\star} = \sqrt{1 - c^{2}}, \qquad V^{2} + (D^{\star})^{2} = c^{2} + (1 - c^{2}) = 1, \]
which is the conceptual core of Eq.~\eqref{eq:circle}: the visibility is the marker overlap itself, because interference between the two augmented paths is weighted by how much their marker routes overlap, and the distinguishability is what remains of unity once that overlap is squared away, because two states of overlap $c$ can be discriminated with a maximal reliability that standard discrimination theory fixes at $\sqrt{1 - c^{2}}$. The two ingredients carry different statuses and the audit at the end of this section keeps them apart: $V = c$ is the kernel's own cross-term overlap; $D^{\star} = \sqrt{1 - c^{2}}$ is imported as the standard optimal-discrimination result, pending the in-formalism definition of the placement supremum (Section~\ref{sec:beyond}). An illustrative check, with the running example $\alpha = 30^\circ$: $c = \cos 60^\circ = \tfrac12$, so $V = \tfrac12$, $D^{\star} = \sqrt{1 - \tfrac14} = \tfrac{\sqrt3}{2}$, and $\tfrac14 + \tfrac34 = 1$; this is a check on the reader's understanding, not a result. The analyzer, finally, does not create distinguishability: it chooses a direction on the marker sphere along which the existing difference between the two marker states is read out, one direction (the bisector between the tags) reading out nothing and restoring full fringes within its port, the perpendicular direction on the sphere reading out the most---and the optimum is not the original tag axis, which reads out strictly less. That is why $D^{\star}$ must be defined as an optimum over admissible record placements rather than read off the obvious detector basis.



While the which-tag attribute stays blank in principle, the two tagged routes stand as a single agreeing-regime superposition. The analyzer's pass-or-block outcome is a record---a further determinate node on the sphere---and here the disagreeing regime enters, with it the first place weights rather than amplitudes compose: the structure is pruned there into the two ports as a disagreeing superposition whose weights are the readout probabilities,
\begin{equation}\label{eq:ports}
p_\theta \;=\; \tfrac{1}{2}\big(1 + c \cos 2\theta\big),
\end{equation}
(three quarters and one quarter at $\alpha = 30^\circ$, analyzer vertical). This split (Eq.~\eqref{eq:ports}) is both regimes standing in one experiment: a disagreeing record layer, the two ports as pockets whose weights add, each pocket carrying a surviving agreeing interior whose amplitudes add into that port's fringes or anti-fringes. Nothing is sorted by hand; the port record partitions, and the interiors interfere within the partition. The partition can be exhibited as a reduced-form calculation rather than described. Let the candidate family be the four augmented routes $T = \{p_L^{+}, p_R^{+}, p_L^{-}, p_R^{-}\}$, with $p_j^{\pm}$ the route through slit $j$ compatible with analyzer record $\pm$. The record induces the two reduced forms
\begin{equation}\label{eq:ledger}
R_{+}(T) = \{p_L^{+}, p_R^{+}\}, \qquad R_{-}(T) = \{p_L^{-}, p_R^{-}\},
\end{equation}
which are mutually exclusive at the record level, $R_{+}(T) \perp R_{-}(T)$, since $+$ and $-$ are exclusive records; while within each, $p_L^{\pm} \not\perp p_R^{\pm}$ at the path level, since no port records which slit. Hence
\[
\boxed{\text{weights compose between } R_{+} \text{ and } R_{-}; \quad \text{kernel amplitudes compose within each } R_{\pm}.}
\]
The specifically Aethic content of the construction is this chain---prune by record, into disagreeing pockets carrying weights, each containing an agreeing interior carrying amplitudes---and the trigonometry below is what the kernel supplies once the chain is fixed. 

(Here the agreeing regime survives inside the exclusive one, and amplitudes compose within each port: the coherent case differs from a classical mixture exactly because the marker states remain one agreeing cluster rather than record-separated components.) Within each port the agreeing pair survives --- passing did not resolve the which-tag attribute, both tags projecting onto the port with nonzero amplitude --- so the pair stands tilted: with $c_1 = \cos(\theta - \alpha)$ and $c_2 = \cos(\theta + \alpha)$, the routes combine at amplitude grade. The tilt is not applied by the analyzer. The coefficients are the kernel contributions the surviving augmented routes already carry, $\mathcal{K}(p_L^{+}) = c_1\, e^{iS_L/\hbar}$ and $\mathcal{K}(p_R^{+}) = c_2\, e^{iS_R/\hbar}$, the analyzer being part of the path structure and each full route through it carrying the corresponding overlap. Conditioning on the port restricts the cluster, $T \to T \cap E_{+}$, and computes $W(T \cap E_{+})$ with the same kernel; renormalization then divides every surviving contribution by one common factor and cannot alter the ratio $c_1 : c_2$. No surviving amplitude is edited; only the domain of surviving paths changes. So the port pattern is
\begin{equation}\label{eq:portamp}
\mathcal{A}_{+}(x) \;\propto\; c_1\, e^{iS_L/\hbar} + c_2\, e^{iS_R/\hbar},
\end{equation}
and squaring it, with $k\,\Delta(x) = (S_L - S_R)/\hbar$ the spacetime phase difference,
\begin{equation}\label{eq:twotag}
I_\theta(x) \;\propto\; c_1^2 + c_2^2 + 2\, c_1 c_2 \cos\!\big(k\,\Delta(x)\big), \qquad\quad V_\theta \;=\; \frac{2\,|c_1 c_2|}{c_1^2 + c_2^2},
\end{equation}
full fringes at the erasing angle, contrast falling to $c$ at the Bloch-orthogonal straddle ($45^\circ$ on the bench, the sphere doubling physical angles), anti-fringes at the complement --- the contrast a pure Bloch-face ratio, the stripe geometry pure spacetime. And the surviving pair is no bookkeeping ghost: its screen face is the erased-port conditional pattern, identically in the analyzer angle --- one residual object under two cuts, the framework's surviving pair and the laboratory's sorted statistics the same thing. \emph{Across} the ports, weights add, and the step is worth writing out. The complementary port carries $d_1 = -\sin(\theta - \alpha)$ and $d_2 = -\sin(\theta + \alpha)$, so
\begin{equation}\label{eq:nosig}
c_1 c_2 + d_1 d_2 \;=\; \cos(\theta - \alpha)\cos(\theta + \alpha) + \sin(\theta - \alpha)\sin(\theta + \alpha) \;=\; \cos 2\alpha,
\end{equation}
and, with $c_j^{2} + d_j^{2} = 1$, $I_{+}(x) + I_{-}(x) \propto 2 + 2\cos 2\alpha\, \cos\!\big(k\,\Delta(x)\big)$, with no $\theta$ anywhere: the unconditional pattern has visibility $c$ at every analyzer angle --- \emph{no-signaling arriving as an algebraic identity of the two composition laws rather than as a constraint checked afterwards} --- a \emph{placement} result, definitional once the composition laws are fixed, and one every interpretation obtains from linearity by some route; what is claimed is the route's directness, not a novel consequence. The record's separating capacity completes the accounting: the two tags answer the analyzer with probabilities $c_1^2$ and $c_2^2$, their ensemble gap the analyzer-indexed distinguishability $D(\theta) = |\sin 2\theta|\,\sin 2\alpha$ --- a record quantity, never a screen quantity, and distinct from the port-normalized $P_\theta$ below; it is $D$ itself at the analyzer placement, and its supremum over the analyzer family, $D^{\star}$, sits at the straddle and equals $\sin 2\alpha = \sqrt{1 - c^{2}}$ --- exactly the Englert distinguishability for pure markers --- so $D^{\star}$ is here exhibited as an in-formalism supremum over the analyzer family, and the corresponding open construction narrows, via the route developed in the long paper \cite{Aethus} and summarized in Section~\ref{sec:beyond}, to two premises: (R1) that an admissible record placement is a cell family over the record register, and (R2) the identification of the record register's state space with the marker sector. (The unconditional pair accordingly saturates the circle exactly at the straddle: $c^{2} + \sin^{2} 2\alpha = 1$.) The linear split $c_D = D(\theta)$, $c_A = 1 - D(\theta)$ is a fact about the \emph{record statistics}, not about the state: by the coupling lemma the two outcome distributions at the analyzer placement $\theta$ decompose into a common part of weight $1 - D(\theta)$ and disjoint parts of weight $D(\theta)$.

The visibility and predictability expressions are now two components of one normalization identity: with $V_\theta = 2|c_1 c_2|/(c_1^2 + c_2^2)$ and the port \emph{predictability} $P_\theta = |c_1^2 - c_2^2|/(c_1^2 + c_2^2)$---predictability rather than distinguishability, since after the analyzer no marker remains and this is Greenberger and Yasin's $P$ \cite{greenberger1988}, the identity below being their $V^2 + P^2 = 1$ for a pure two-path state---
\begin{equation}\label{eq:portidentity}
(2 c_1 c_2)^{2} + (c_1^{2} - c_2^{2})^{2} \;=\; (c_1^{2} + c_2^{2})^{2} \qquad\Longrightarrow\qquad V_\theta^{2} + P_\theta^{2} \;=\; 1,
\end{equation}
the interference term standing as $V_\theta$ and the predictability term as $P_\theta$, the two complementary pieces of one normalization, so that the quadratic relation is not another empirical input --- a port-wise identity of the construction that has the circle's shape while involving $P_\theta$ and not $D$, and a second exact circle beside the roadmap's unconditional one. Two relations are in play, not one read twice: at weight grade the record sector splits linearly, $c_D = D_p$ and $c_A = 1 - D_p$ at each placement, (the total-variation decomposition of that placement's marker statistics), while the square map joins the two faces at the optimal placement by relating the kernel's overlap to the record capacity, $c^2 + (D^{\star})^2 = 1$. One configuration, both composition laws, every factor assigned to its face.







\subsection{Beyond the Circle: Strictness, Leakage, and the Optimality Frontier}\label{sec:beyond}

The circle of Eq.~\eqref{eq:circle} is an equality for pure markers, and the two ways below it, a leak coupled to the path and a leak confined to the marker sector, move the point strictly inside it by the amount the coupling lemma fixes; the limiting consistency checks at the two nodes, the in-formalism definition of an arbitrary record placement as a cell family over the record register, and the optimality frontier of record placements are developed in the long paper \cite{Aethus}, and nothing below depends on them.

\subsection{The Bell Leg: Realism Without a Hidden Valuation}\label{sec:bell}

% >>> ADDED (Fable 5.1, Sep 2 2026, author-directed re-centering): the Deflation Reductio ported from the
% companion into neutral vocabulary as the third exhibit --- what the recovery's realism is, shown against
% Bell's theorem.
The recovery delivers a realism of a specific shape: enough counterfactual structure to say what an experiment would have shown, and no hidden valuation for Bell's theorem to kill. The second half is what a reader will doubt, since the second postulate's reading of a blank as \emph{open} is what the deflationary reading --- a blank is an unknown value, determinate beneath --- denies, and the two readings agree on every classical statistic. The exhibit is a reductio showing where they part, and that they part exactly where the statistics violate Bell's bounds.

\begin{theorem}[Deflation reductio]\label{thm:deflation}
Call an assignment $\delta$ a \emph{lossless deflation} of an information state $A$ when it provides, for $A$ and every state $B$ in $A$'s refinement structure, a family $W(B)$ of complete valuations of $B$'s attributes such that: (D1) each $w \in W(B)$ assigns every attribute one value, independent of which other attributes are stated or measured, agreeing with $B$'s stated entries; (D2) on the exclusive-regime fragment --- states all of whose cross content is declared --- $B$ is valid iff $W(B) \neq \varnothing$, with refinement restricting the family; (D3) the construction's ensemble statistics over any attributes of $A$ are the marginals of one probability measure $\mu$ on $W(A)$, each run's outcomes identified with that run's values, the choice of which attributes are measured being independent of $w$. Then: (a) if the blank attributes of $A$ contain a Bell quadruple --- four attributes whose CHSH combination exceeds $2$ under the kernel's own weights or under the observed statistics --- no lossless deflation of $A$ exists; (b) if every correlation among the blank attributes satisfies the Bell inequalities of its marginal scenario---for a Bell quadruple the CHSH bounds, which are necessary and sufficient there by Fine's theorem, and in general the facet inequalities of the corresponding correlation polytope \cite{pitowsky1994}---the obstruction of~(a) is absent: Fine's theorem gives a joint distribution for the specified marginals, which removes that obstruction without establishing a lossless deflation, since a deflation needs valuations compatible across the whole refinement structure and the marginal theorem supplies none.
\end{theorem}
\noindent\textit{Proof.} (a) Suppose $\delta$ exists. By (D1) each $w \in W(A)$ assigns a value to all four attributes, so the $\mu$ of (D3) is a measure on the four-fold value space, and by (D3) the four pairwise correlations are its marginals, measurement independence letting one $\mu$ serve all four setting pairs. By Fine's theorem \cite{fine1982} the existence of such a joint distribution is equivalent to the CHSH inequalities \cite{pitowsky1994}, so the combination is at most $2$, contradicting the hypothesis. (b) Bell-bounded correlations admit a joint distribution with the given marginals, Fine's theorem in its converse direction, and Pitowsky's correlation polytope in general; that is all (b) asserts. \hfill$\square$

\noindent What the theorem fixes is exact: to read a blank as an unknown in the deflationary way is to posit a lossless, noncontextual, measurement-independent assignment, and that assignment's image is a single sample space the Bell statistics exclude; a deflation that translates openness rather than preserving it, and a contextual assignment, lie outside the theorem's class and are addressed in the determinacy companion \cite{AethusRealist}.




\paragraph{A worked instrument: the bomb tester.} The Elitzur--Vaidman arrangement \cite{elitzur1993quantum} exercises the formalism on a trace-decreasing survival operation and its complementary absorption event. With arms $u$ (safe) and $v$ (sensor), the second splitter sending $|u\rangle \mapsto (|b\rangle + |d\rangle)/\sqrt 2$ and $|v\rangle \mapsto (|b\rangle - |d\rangle)/\sqrt 2$, a transparent dud gives unit bright-port weight, while an ideal live absorber sends the arm-$v$ amplitude to an orthogonal absorbed state and leaves $\tfrac12|b, L\rangle + \tfrac12|d, L\rangle + \tfrac{1}{\sqrt 2}|E\rangle$: explosion $\tfrac12$, bright $\tfrac14$, dark $\tfrac14$. The live device, taken as live, realizes a record-compatible restriction on the path space with record partition $\{v, u\}$, the survival operation $K_{\mathrm{survive}} = |u\rangle\langle u|$ being its projector, so conditioning on survival is the L\"{u}ders conditioning of clause (v) on the enlarged system of photon, bomb and record register, within the projective class the theorem treats; that fixed-live arm partition is distinct from the live/dud status partition, which the long paper separates by a persisting record stating the live response. What the classification adds is separate: relative to a state stating the live response, the candidate content of two surviving amplitudes recombining is jointly invalid with the explosion attribute at either value, and Lemma~\ref{lem:record-exclusion-q} excludes it when the explosion and intactness records are persistent, a premise about a macroscopic bomb; the mediating attribute is the bomb's response at the trial event, not a record of the photon's path. The weights are the kernel's and the response law's, the classification fixes the composition law they enter under, and the coherent-device case, decided in the long paper on the definition of the would-clause and a detector premise as a proposed divergence from universal unitarity, is not decided by this example.



\section{Placement, Costs, and Open Constructions}\label{sec:discussion}\label{sec:costs}

\emph{Placement.} Against the operational reconstructions \cite{hardy2001, cbh2003, cdp2011, masanesmuller2011}, which take measurement and a probabilistic structure as primitive and derive the Hilbert form, the present base is a logic of indexed information states from which measurement structure is organized by records and the Hilbert form is a representation on one finite domain. Against relational quantum mechanics and QBism it shares the indexing of facts and adds an ontology for the superposition and a derived rather than posited criterion for when a cluster is admissible; against the dynamical-collapse programs it shares only the position that coherence ends somewhere for observer-grade systems, with no mechanism and no rate, the criterion being the third postulate under the record premise; and determining the signature partition in realistic configurations is the dynamical work decoherence theory does, which this paper takes as input. The comparisons are developed in the companions \cite{Aethus, AethusBrief}, and nothing in the theorems above turns on them.





\emph{The nested-observer protocols, in outline.} The $\psi$-ontology constraint is a different structure, one preparation device ranging over nonorthogonal pure states, and its treatment stays with the companion. The Frauchiger--Renner protocol's load-bearing premise is that one agent holds a coherent state of a colleague's whole laboratory across the colleague's outcome alternatives, and that coherence claim is what the framework assesses under the record-admissibility criterion. The inference that fails if it is excluded is the certainty step, that a superobserver's ok outcome fixes a friend's outcome. The long paper computes the four-outcome distributions over $(\mathrm{ok}_1, \mathrm{ok}_2)$, $(\mathrm{ok}_1, \mathrm{fail}_2)$, $(\mathrm{fail}_1, \mathrm{ok}_2)$ and $(\mathrm{fail}_1, \mathrm{fail}_2)$: $(\tfrac{1}{12}, \tfrac{1}{12}, \tfrac{1}{12}, \tfrac34)$ for the coherent assignment $(|h,{\downarrow}\rangle + |t,{\downarrow}\rangle + |t,{\uparrow}\rangle)/\sqrt 3$, and $(\tfrac14, \tfrac14, \tfrac14, \tfrac14)$ for the mixture over its three pockets at weight $\tfrac13$ each \cite{Aethus}. The detector and the friend instantiate one formal criterion, and neither exclusion requires that the external observer has read the record: present sealing alone does not establish that the record-access refinements on which the third postulate operates have been removed.
\begin{lemma}[Conditional record exclusion]\label{lem:record-exclusion-q}
Let $A_{\mathrm{coh}}$ be a candidate state asserting a specified maintained coherence, with declared descent reflexive and transitive and validity read base-generated, and let the relevant base contradictions be the joint statings of that coherence with an incompatible separating record. If the record family is persistent over the candidate's declared descent, every declared refinement of $A_{\mathrm{coh}}$ having a further descendant stating a separating record, then $A_{\mathrm{coh}}$ is invalid under $P_3$.
\end{lemma}
\begin{proof}
Validity requires a safe refinement, all of whose declared descendants lie outside the base; by persistence every candidate safe refinement has a base-contradictory descendant, so none exists.
\end{proof}
\noindent Establishing the lemma's premise at the scope of the whole laboratory, the composite proposed as coherently controlled, is the outstanding physical question; a path subsystem may lack interference while the system--marker composite retains coherence, so the records that matter are those of the whole laboratory. A physical realization must specify the candidate configurations, their declared refinements and the record incompatibilities on which $P_3$ operates. Under the record premise it is the coherent preparation that is excluded, not the protocol's measurement: applied to the disagreeing mixture the framework assigns, the measurement yields the mixture's statistics, and whether the superposition-basis operation on an observer-containing laboratory can be implemented is left to a physical realization. The argument consumes the record--coherence incompatibility, stated in full in the companion \cite{Aethus}, and its trigger, in-principle irretrievability, has exact conditions that remain open \cite{AethusBrief}; a perfectly sealed laboratory is where the trigger is contested, and the verdict is conditional accordingly.

\emph{Costs, inventoried.} The assumptions are the package $\mathsf{Q}_{\log}$, $\mathsf{Q}_{\mathrm{ker}}$, $\mathsf{Q}_{\mathrm{phys}}$ and $\mathsf{Restr}$ of Section~\ref{sec:recovery}, itemized per result in Table~\ref{tab:kernel-boundary}: the weights posited by the kernel, the bridge to physical unrealizability a principle, the restriction stipulated, and the observer applications' record premise named rather than derived.

\emph{Open constructions, as determinate targets.} A minimality theorem for the pruning rule---which assumptions on the rule class are necessary and sufficient for the $\exists\forall$ form, with an explicit countermodel for each assumption dropped, which would turn ``selected within a class'' into a claim that cannot be dismissed as framework-dependent; the measure-invariance proof for the reopened-blank measure; the in-formalism definition of the distinguishability supremum (narrowed by the two-tag construction's analyzer-family supremum, and now further by the conditional route of the long paper, summarized in Section~\ref{sec:beyond}, to the two named premises (R1), the cell-family definition of an admissible record placement, and (R2), the state-space identification of the record register, with the state-space identification (R2) the next target); the kernel's derivation; the exact trigger conditions for in-principle irretrievability, and, added here, the continuum extension of the classification (almost-identical signatures in place of identical ones, ``permanently indistinguishable'' becoming a limit notion, which is the approximateness charged to decoherence) --- on this one flank the framework shares with the dynamical-collapse programs the position that coherence ends somewhere for observer-grade systems, and nothing of their mechanism: no criterion is posited beyond the third postulate, which excludes a candidate when every declared refinement of it retains a contradiction-reaching descendant, the same rule that ends interference when a detector carries a distinguishing record, so no cut is drawn and no rate or scale enters; what the framework stakes is one physical claim, that an observer-grade record structure cannot be kept from exporting its records, and that claim is its one structural divergence from universal unitarity --- two disciplines keeping it from going verbal: ``unblockable'' must not be cashed as what unitary dynamics already calls irreversible decoherence, on which reading the two sides agree on every experiment and the claim is a gloss, so the discriminating regime lies among records scrambled but present, the spin-echo and Loschmidt-echo regime; and the interferometry program climbing toward ever-richer record-bearing systems held coherent can only fail to find the point at which the escape closes, every successful recombination reading as ``records still blockable'' until the conditions are named --- so the stake is a commitment to a prediction rather than a prediction, priced as exactly that here and in the overview \cite{AethusBrief}; the record-individuation lemma; and the extension of the node-grade construction to general spin dynamics, higher-spin systems, and further internal arenas. Each is a mathematical question with a stateable success condition, and the framework binds itself to falsifiable posture where it can --- most sharply at the negativity--contextuality joint, where the refusal of a single sample space is a standing structural commitment rather than a preference, (the compliance riding on that refusal and not on the signature of the tree's weights, which the fuller development admits signed while keeping the term algebra inverse-free \cite{Aethus}).

\emph{The close, in one paragraph.} Under a complete and persistent record structure the third postulate classifies the subsets of a path space, and given the kernel the classification's clusters are represented in a Hilbert space with the Born statistics, record observables, and a stipulated restriction realized as projection, all on one finite domain. The measurement rule is relocated to that restriction, not eliminated, and the observer applications rest on a physical record premise stated as such.

\vspace{0.4cm}
\noindent{\small \textbf{Provenance note.} The framework, its postulates, and their applications are the author's, established and outlined end-to-end in the author's own terminology by May 2023 and deposited at PhilArchive from November 2024 onward, with successive deposits dating the later-coined terms ahead of any AI involvement; the continuity is laid out deposit-by-deposit in \cite{Aethus}. AI involvement, stated once and consistently: the framework, the postulates, and the results deposited through 2025 are the author's; from June 2026 an AI collaborator drafted derivations and translated outlined structures into the standard formalism within the collaboration, and the delta of this draft beyond the deposited version---the triangle and monotone re-classification theorems, the leakage-geometry analysis, and the compressed conditional route---is of that mixed provenance, corrected across several AI referee rounds, all within one model family; a dated section-level audit travels in the companion session record. Independent human refereeing is the gate before any submission, and AI review of any kind is not that gate. The disclosure conventions the companion documents in full apply here.}
% TO DO before submission: check the target journal's AI-use policy against this note, which is written to satisfy a drafting-assistance disclosure and to exclude AI authorship.

\subsection{A finite correspondence model}\label{sec:correspondence}

The compatibility question can be made exact on one finite model. Take a balanced Mach--Zehnder interferometer with arms $u$ and $d$ and output ports $b$ and $k$, the kernel assigning the four routes $\mathcal{K}(u \to b) = \mathcal{K}(d \to b) = \mathcal{K}(u \to k) = \tfrac12$ and $\mathcal{K}(d \to k) = -\tfrac12$, so that the prepared cluster $T_0$ of all four routes has $\Phi(T_0) = |b\rangle$. The carriers are the post-record state $A$, in which the port record is declared and the port is blank, and its two record cells, $B$ with the port stated as $b$ and $K$ with it stated as $k$; the declared decomposition is $A \rightsquigarrow \gamma_b B + \gamma_k K$; and the kernel weights are $W(B) = |\tfrac12 + \tfrac12|^2 = 1$ and $W(K) = |\tfrac12 - \tfrac12|^2 = 0$, $K$ being the dark port. Two readings of the structural coefficients are compared: the present model's, (S), every record cell at a positive coefficient, $\gamma_b, \gamma_k > 0$; and the kernel-tied reading, (W), $\gamma_j = W_j$, under which $K$ is no proper child of $A$. With no base contradiction both readings make $A$ valid and agree on every validity judgment; they differ only in the proper-child relation, which contains $K \le A$ under (S) and not under (W). Add one base contradiction, declaring the bright-port cell contradicted, $B \in \mathrm{Base}$. As a formal countermodel this is enough to test the two relations; a physical realization of it, a record that the photon did not leave port $b$, would require specifying the enlarged state, its record consistency, and whether the original preparation and kernel still apply, none of which the model does. Under (S), $A$ is valid, its proper child $K$ having a cone free of base contradictions; under (W), $A$ is invalid, its only proper child of positive weight being contradicted. The divergence is an instance of a general fact about finite models.

\begin{proposition}[Where the two readings part]\label{prop:readings}
Let a finite model's declared descent consist of record cells whose kernel weights are nonnegative and add over their sub-cells, and let $\le_S$ and $\le_W$ be its proper-descendant relations under (S) and (W), $\le_W$ omitting every cell of zero kernel weight. For every state $A$ of positive weight, validity under (W) implies validity under (S); and $A$ is valid under (S) but not under (W) exactly when $A \notin \mathrm{Base}$, every proper descendant of $A$ of positive weight has a base contradiction in its $\le_W$-cone, and some proper descendant of zero weight has none in its $\le_S$-cone.
\end{proposition}
\begin{proof}
By Proposition~\ref{prop:closes}, validity is $A \notin \mathrm{Base}$ together with a proper descendant whose cone avoids $\mathrm{Base}$. Weights being nonnegative and additive, a sub-cell weighs no more than its cell, so every proper descendant of positive weight is reached through cells of positive weight, and $\le_S$ and $\le_W$ agree on such descendants. If $B \le_W A$ has a $\le_W$-cone free of base contradictions, follow cells of positive weight down from $B$ to a cell with no sub-cells, which exists because the model is finite and a cell of positive weight with sub-cells has one of positive weight; that cell lies in $B$'s $\le_W$-cone, so it is not in $\mathrm{Base}$, and its $\le_S$-cone is itself, so it witnesses validity under (S). Conversely, an (S)-witness of positive weight is a proper descendant under (W) whose $\le_W$-cone lies inside its $\le_S$-cone, so it witnesses validity under (W). The readings therefore differ only where (S) is witnessed by a descendant of zero weight and (W) by none, which is the stated condition. A check over the 23{,}478 states of positive weight in 4{,}000 randomly generated finite models, made by the accompanying script \texttt{readings\_check.py} (run as \texttt{python3 readings\_check.py 4000 20261006}), agrees in every case.
\end{proof}

\noindent The two readings therefore agree on every validity judgment but one kind: whether a state is valid when its only continuations free of contradiction have kernel weight zero. Under (S) such a state is valid, carried by an outcome of zero probability on which conditioning is undefined; under (W) it is invalid. Within this finite additive record-cell model, the validity comparison reduces to whether kernel-null descendants may supply a safe witness. The model does not make that decision, and it does not reach the framework's other relative-negligibility cases, components carried in different measure dimensions among them; which reading the physics requires is left open.

\printbibliography

\end{document}